% Le titre du chapitre et son label sont declares dans main.tex
% (cf. \chapter{Materiel et methodologie} avant % Le titre du chapitre et son label sont declares dans main.tex
% (cf. \chapter{Materiel et methodologie} avant % Le titre du chapitre et son label sont declares dans main.tex
% (cf. \chapter{Materiel et methodologie} avant % Le titre du chapitre et son label sont declares dans main.tex
% (cf. \chapter{Materiel et methodologie} avant \input{chapters/chapter2.tex}).

\phantomsection
\section*{Introduction du chapitre}

Après avoir posé le cadre conceptuel et exposé l'état de l'art, nous présentons le matériel mobilisé et la démarche méthodologique suivie. Ce chapitre fait le lien entre la réflexion théorique et la mise en œuvre technique.

Nous présentons le matériel mobilisé, notamment le corpus documentaire, les ontologies candidates et l'environnement technique. Nous décrivons les huit phases de la démarche suivie, de l'étude comparative des ontologies à l'évaluation du système, puis en proposons une synthèse.

% ============================================================
% II.1. MATÉRIEL
% ============================================================

\section{Matériel}
\label{sec:materiel}

Le matériel comprend le corpus documentaire, les ontologies candidates et l'environnement technique. Les deux tableaux qui suivent récapitulent ces ressources, réparties entre documentation, modèles conceptuels et outils techniques.

% ------------------------------------------------------------
% II.1.1. Corpus documentaire
% ------------------------------------------------------------

\subsection{Corpus documentaire}
\label{subsec:corpus}

Nous avons choisi l'ontologie avant de constituer le corpus afin de déterminer les notions et les relations qui guideraient l'extraction. Nous avons réuni les publications scientifiques et les données institutionnelles consultées sur burkinaherbiermed.org et auprès de l'INERA.

La figure \ref{fig:flux_corpus} résume le processus de constitution du corpus, du choix de l'ontologie à la sélection des documents.

\begin{figure}[H]
\centering
\includegraphics[width=0.8\textwidth]{flux_corpus.png}
\caption{Constitution du corpus documentaire}
\label{fig:flux_corpus}
\end{figure}

\subsubsection{Sources scientifiques publiées}

Notre corpus scientifique comprend 35 rapports de recherche, articles, thèses et mémoires provenant de revues, d'universités et d'institutions de recherche burkinabè. Ces travaux ethnobotaniques, ethnopharmacologiques et phytochimiques documentent les noms scientifiques et vernaculaires, les parties utilisées, les pathologies traitées, les modes de préparation et, pour les études qui les rapportent, les dosages et les données de toxicité.

Nous avons retenu ces travaux pour leur qualité et leur pertinence au contexte burkinabè, des publications fondatrices aux plus récentes. Cette diversité d'espèces, de pathologies et de pratiques thérapeutiques soutient l'extraction. Quatre travaux illustrent les types de données mobilisés dans le corpus :
\begin{itemize}
\item Zerbo et al.~\cite{zerbo_plantes_2011} ont étudié les plantes médicinales chez les Sanan du Burkina Faso. Leur travail documente 85 espèces utilisées pour traiter 63 affections.
\item Ouédraogo et al.~\cite{ouedraogo_revue_2024} ont dressé une revue des plantes médicinales à potentiel aromatique du pays, recensant 42 espèces et en caractérisant les propriétés et les usages.
\item Tibiri et al.~\cite{tibiri2020} ont mené une enquête à l'échelle nationale sur les plantes utilisées contre les maladies hépatiques, documentant plus de 2 000 entrées de plantes correspondant à 212 espèces.
\item Ouédraogo et al.~\cite{ouedraogo2020markets} ont enquêté auprès de 30 commerçants de plantes médicinales dans six grandes villes, identifiant 104 espèces utilisées pour diverses pathologies.
\end{itemize}

Ces travaux fournissent des méthodes réutilisables et des observations de terrain recueillies au Burkina Faso.

\subsubsection{Données de l'Herbier National des Plantes Médicinales du Burkina Faso}

Le site burkinaherbiermed.org diffuse des données sur les plantes médicinales du pays, notamment des informations botaniques, des usages thérapeutiques et des références bibliographiques. Ces données proviennent du projet « Herbier avec chimiothèque », piloté par le Ministère de la Santé. Ce projet a pour objectif de constituer un herbier national dédié aux plantes médicinales et une chimiothèque de marqueurs chimiques. Il documente les plantes utilisées contre les dix pathologies prioritaires définies par le ministère, géolocalise les stations et conserve des spécimens. Ces informations alimentent une base nationale accessible par application mobile et site web. Le projet associe un herbier, une chimiothèque et une base de données ethnobotaniques et géographiques. Il complète les autres herbiers burkinabè, notamment le HNBU du CNRST/INERA, le Centre National de Semences Forestières et l'herbier de l'Université Joseph Ki-Zerbo.

\subsubsection{Données de l'INERA}

La base de données de l'Institut de l'Environnement et de Recherches Agricoles (INERA) comprend plus de 7 000 espèces végétales et les coordonnées géographiques des stations où elles ont été observées. Ces données complètent les informations ethnobotaniques par une dimension spatiale et permettent d'étudier la répartition des plantes au Burkina Faso. Le rapprochement des usages rapportés et de la distribution géographique peut faire apparaître des régularités régionales, qui restent à vérifier.

Le tableau \ref{tab:corpus} récapitule les sources du corpus.

\begin{table}[H]
\centering
\caption{Sources du corpus documentaire}
\label{tab:corpus}
\begin{tabular}{|L{4cm}|L{2.8cm}|L{3cm}|L{5cm}|}
\hline
\textbf{Source} & \textbf{Type} & \textbf{Volume estimé} & \textbf{Contenu principal} \\
\hline
Articles et thèses & Publications scientifiques & 35 documents & Études ethnobotaniques, pharmacologiques \\
\hline
burkinaherbiermed.org & Base de données en ligne & Plusieurs centaines de plantes & Données botaniques, usages, références \\
\hline
INERA & Base de données institutionnelle & > 7 000 espèces & Coordonnées géographiques, identification \\
\hline
\end{tabular}
\end{table}

Les documents du corpus relèvent de types d'étude différents, et chacun est exploitable pour les informations qu'il rapporte. Trois profils se distinguent. Les enquêtes ethnobotaniques et ethnopharmacologiques, menées auprès de tradipraticiens ou dans les marchés, rapportent les noms vernaculaires, les parties utilisées, les modes de préparation et les pathologies traitées, souvent plante par plante. Les études expérimentales, notamment les essais \textit{in vitro} et les travaux phytochimiques, mesurent des activités pharmacologiques, des concentrations, des niveaux de preuve et parfois les lieux de récolte. Elles renseignent moins sur les noms vernaculaires et les pratiques de préparation. Enfin, les revues de synthèse dressent des inventaires à grande échelle et servent surtout à repérer les sources primaires.

Nous avons appliqué trois règles à l'extraction. Chaque source ne renseignait que les champs qu'elle mesurait. Nous avons laissé vides les champs non couverts, sans les compléter par déduction. Nous n'avons pas retenu comme fait établi un usage traditionnel rapporté dans un article expérimental sans référence ethnobotanique. Enfin, nous n'avons pas additionné les nombres d'espèces d'une source à l'autre, car une revue peut reprendre les espèces documentées dans les enquêtes qu'elle cite. Le tableau \ref{tab:corpus} présente les sources document par document, sans agréger les volumes.

Ce corpus rassemble des études menées dans plusieurs régions du Burkina Faso, portant sur différentes pathologies et mobilisant des approches méthodologiques variées. Il a servi de base à l'extraction des connaissances et à l'enrichissement de l'ontologie. Le choix préalable de l'ontologie a guidé l'extraction et l'organisation des documents selon les concepts et les relations définis.

% ------------------------------------------------------------
% II.1.2. Ontologies candidates
% ------------------------------------------------------------

\subsection{Ontologies candidates}
\label{subsec:ontologies_candidates}

Nous avons retenu cinq ontologies candidates pour l'étude comparative. Le chapitre \ref{chap:cadre_theorique}, aux sections \ref{subsubsec:ontologies_internationales} et \ref{subsubsec:ontologies_africaines}, les présente. Elles couvrent différents domaines, notamment la botanique, la médecine traditionnelle, la pharmacologie et l'ethnomédecine. Le tableau \ref{tab:ontologies_candidates} en présente les caractéristiques.

Nous les avons comparées à l'aide de la grille de critères présentée en phase 1 (section \ref{subsec:phase1}). Nous avons sélectionné l'ontologie la plus adaptée à notre contexte avant de constituer le corpus, afin de fixer le modèle conceptuel qui a guidé l'extraction.

\begin{table}[H]
\centering
\caption{Ontologies candidates pour la comparaison}
\label{tab:ontologies_candidates}
\begin{tabular}{|C{0.5cm}|L{3cm}|L{3cm}|L{3cm}|L{2cm}|L{2.5cm}|}
\hline
\textbf{N°} & \textbf{Ontologie} & \textbf{Domaine} & \textbf{Zone géographique} & \textbf{Format} & \textbf{Source d'accès} \\
\hline
1 & WATRIMed & Médecine traditionnelle & Afrique de l'Ouest & OWL, graphe de connaissances & IOS Press (2 articles) \\
\hline
2 & ontoMEDTRAD & Médecine traditionnelle & Afrique de l'Ouest & OWL + iconique & Thèse A. Kouamé (UGB) \\
\hline
3 & ATHMed & Médecine traditionnelle & Afrique (Ghana) & Cadre conceptuel & Frontiers in AI \\
\hline
4 & Deep Ontologies & Médecine traditionnelle & Afrique & OWL (Protégé) & Articles ARIMA \\
\hline
5 & FLOPO, PO & Botanique & International & OWL & OBO Foundry \\
\hline
\end{tabular}
\end{table}

Chacune présente des acquis et des lacunes. Les ontologies botaniques (FLOPO, PO) décrivent bien la plante mais ne traitent pas des usages médicinaux. Les ontologies pharmaceutiques (DrOn) se resserrent sur les médicaments modernes. L'OTM, prometteuse, reste sans ancrage africain. Les initiatives africaines (WATRIMed, ontoMEDTRAD, ATHMed) s'intéressent en revanche à la médecine traditionnelle du terrain, avec ses dimensions culturelles et ses recettes multi-plantes, même si ATHMed en reste, elle, au stade d'un cadre proposé. La phase 1 en livre la comparaison détaillée.

% ------------------------------------------------------------
% II.1.3. Outils et environnement technique
% ------------------------------------------------------------

\subsection{Outils et environnement technique}
\label{subsec:environnement}

Les outils techniques occupent une place importante dans notre démarche. Nous décrivons les phases concernées et les raisons de leur choix.

La figure \ref{fig:dispositif_technique} résume l'environnement d'expérimentation, en précisant où chaque composant est exécuté et comment ils communiquent.

\begin{figure}[H]
\centering
\includegraphics[width=0.8\textwidth]{dispositif_technique.png}
\caption{Dispositif technique d'expérimentation}
\label{fig:dispositif_technique}
\end{figure}

\subsubsection{Modèles de langage}

Nous avons utilisé deux modèles, chacun pour une étape distincte. Pour l'extraction des connaissances (phase 5), nous avons utilisé \textbf{gpt-6-luna} d'OpenAI, accessible par l'environnement OpenWork. Pour la génération des réponses du système RAG (phase 7), nous avons utilisé \textbf{Llama 3 8B} \cite{meta2024llama3}, un modèle ouvert. Gemini (\textit{gemini-3.5-flash-lite}) nous a servi de point de comparaison lors de l'évaluation.

Le tableau \ref{tab:llama_acces} présente les modes d'accès à ces modèles.

\begin{table}[H]
\centering
\caption{Modèles de langage et modes d'accès}
\label{tab:llama_acces}
\begin{tabular}{|L{4cm}|L{4cm}|L{6cm}|}
\hline
\textbf{Usage} & \textbf{Modèle} & \textbf{Mode d'accès} \\
\hline
Extraction (phase 5) & gpt-6-luna & OpenWork \\
\hline
Génération RAG (phase 7) & Llama 3 8B & Ollama sur Google Colab, API exposée par ngrok \\
\hline
Comparaison (phase 8) & gemini-3.5-flash-lite & API de Google \\
\hline
\end{tabular}
\end{table}

Nous avons retenu Llama 3 pour la génération parce qu'il s'agit d'un modèle ouvert que nous pouvons exécuter sur une instance que nous contrôlons. La version 8B tient dans la mémoire d'un GPU accessible gratuitement sur Colab. Ollama sert le modèle et ngrok expose son API à l'application web, qui reste sur le poste local.

\subsubsection{Plateforme d'exécution : Google Colab}

Nous avons exécuté Llama 3 8B sur \textbf{Google Colab}, une plateforme de notebooks donnant accès à des GPU adaptés à l'inférence.

\textbf{Phases d'utilisation :} phases 7 et 8 pour la génération des réponses et l'évaluation du système RAG.

Cette plateforme nous a donné accès à un GPU sans nécessiter l'acquisition d'une infrastructure locale. La session reste cependant temporaire et l'adresse exposée par ngrok change à chaque lancement, ce qui convient à un prototype mais pas à un service pérenne.

\subsubsection{Base de connaissances : Neo4j Desktop}

Nous avons utilisé \textbf{Neo4j Desktop}, la version locale de Neo4j, pour créer, gérer et interroger la base de connaissances orientée graphe.

\textbf{Phase d'utilisation :} phase 6 pour la construction et l'interrogation de la base.

Le tableau \ref{tab:neo4j} présente les caractéristiques de Neo4j Desktop.

\begin{table}[H]
\centering
\caption{Caractéristiques de Neo4j Desktop}
\label{tab:neo4j}
\begin{tabular}{|L{5cm}|L{8cm}|}
\hline
\textbf{Caractéristique} & \textbf{Valeur} \\
\hline
Type & Base de données orientée graphe \\
\hline
Langage de requête & Cypher \\
\hline
Version & Desktop (locale) \\
\hline
Format d'import & CSV, JSON \\
\hline
Interface & Neo4j Browser, Bloom \\
\hline
\end{tabular}
\end{table}

Nous avons choisi Neo4j pour représenter les relations entre les plantes, les usages, les propriétés et les sources, interroger le graphe avec Cypher et assurer la traçabilité des connaissances. La base permet également d'ajouter des connaissances et de visualiser les liens entre les entités.

\subsubsection{Orchestration RAG : LangChain}

Nous avons utilisé \textbf{LangChain} pour définir les consignes et appeler les modèles de langage de façon uniforme. La recherche dans Neo4j passe par le pilote officiel et par nos propres requêtes.

\textbf{Phase d'utilisation :} phase 7 pour l'accès aux connaissances et le développement web.

Le tableau \ref{tab:langchain} présente les composants LangChain utilisés.

\begin{table}[H]
\centering
\caption{Composants LangChain utilisés}
\label{tab:langchain}
\begin{tabular}{|L{5cm}|L{8cm}|}
\hline
\textbf{Composant} & \textbf{Rôle} \\
\hline
PromptTemplate & Préparation des consignes de génération \\
\hline
ChatOpenAI & Appel de Llama 3 par l'API d'Ollama, compatible avec celle d'OpenAI \\
\hline
ChatGoogleGenerativeAI & Appel de Gemini pour la comparaison \\
\hline
\end{tabular}
\end{table}

Nous avons retenu LangChain parce qu'il permet de changer de modèle sans modifier le reste du système. Nous n'avons pas confié la sélection des données à une chaîne générant du Cypher : le chapitre \ref{chap:conception_developpement} justifie ce choix \cite{lewis2020rag}.

\subsubsection{Interface web : Streamlit}

Nous avons développé l'application web avec \textbf{Streamlit}, un framework Python qui permet de créer des applications web interactives.

\textbf{Phase d'utilisation :} phase 7 pour le développement de l'application web.

L'application propose un chatbot en langage naturel, des recherches par plante, pathologie ou région, des fiches descriptives et l'affichage des sources et des niveaux de preuve. Nous avons également prévu une fonction de visualisation du graphe de connaissances.

Nous avons retenu Streamlit pour sa simplicité d'utilisation, sa gestion des sessions, son intégration avec LangChain et la rapidité de développement qu'il permet. Il facilite aussi le déploiement de l'application.

\subsubsection{Bibliothèques Python}

Nous avons mobilisé plusieurs bibliothèques Python au cours des phases 5, 6 et 7. Le tableau \ref{tab:bibliotheques} en présente la liste et le rôle.

\begin{table}[H]
\centering
\caption{Bibliothèques Python utilisées}
\label{tab:bibliotheques}
\begin{tabular}{|L{5cm}|L{3cm}|L{5cm}|}
\hline
\textbf{Bibliothèque} & \textbf{Version} & \textbf{Rôle} \\
\hline
\texttt{neo4j} & 5.28 & Pilote de connexion à Neo4j \\
\hline
\texttt{langchain}, \texttt{langchain-core} & 0.3 & Consignes et appel des modèles \\
\hline
\texttt{langchain-openai} & 0.3 & Accès à l'API d'Ollama \\
\hline
\texttt{langchain-google-genai} & 2.1 & Accès à Gemini \\
\hline
\texttt{streamlit} & 1.64 & Interface web \\
\hline
\texttt{plotly} & 7.1 & Visualisation du graphe \\
\hline
\texttt{rdflib} & 7.6 & Lecture du fichier Turtle \\
\hline
\texttt{pytest} & 9.1 & Tests logiciels \\
\hline
\texttt{pyngrok} & Sur Colab & Ouverture du tunnel ngrok \\
\hline
\end{tabular}
\end{table}

Ces bibliothèques couvrent les besoins de la recherche, notamment la connexion à la base, l'orchestration du système RAG, le développement de l'interface, le chargement du modèle, la manipulation de l'ontologie et le traitement des données.

% ------------------------------------------------------------
% II.1.4. Tableau récapitulatif du matériel
% ------------------------------------------------------------

\subsection{Tableau récapitulatif du matériel}
\label{subsec:recap_materiel}

Le tableau \ref{tab:recap_materiel} rassemble l'ensemble du matériel mobilisé.

\begin{table}[H]
\centering
\caption{Récapitulatif du matériel}
\label{tab:recap_materiel}
\begin{tabular}{|L{3cm}|L{4cm}|L{6cm}|}
\hline
\textbf{Catégorie} & \textbf{Élément} & \textbf{Spécification} \\
\hline
\multirow{3}{*}{Corpus} & Documents scientifiques & 35 rapports, articles, thèses \\
\cline{2-3}
& Base burkinaherbiermed.org & Plusieurs centaines de plantes \\
\cline{2-3}
& Base INERA & > 7 000 espèces géolocalisées \\
\hline
\multirow{5}{*}{Ontologies} & WATRIMed & OWL, Afrique de l'Ouest \\
\cline{2-3}
& ontoMEDTRAD & OWL + iconique, Afrique de l'Ouest \\
\cline{2-3}
& ATHMed & Cadre proposé, Ghana \\
\cline{2-3}
& Deep Ontologies & OWL, Afrique \\
\cline{2-3}
& FLOPO, PO & OWL, botanique \\
\hline
Modèles de langage & gpt-6-luna ; Llama 3 8B & OpenWork (extraction) ; Ollama sur Colab et ngrok (RAG) \\
\hline
Plateforme & Google Colab & GPU T4 \\
\hline
Base de connaissances & Neo4j Desktop & Orientée graphe, Cypher \\
\hline
Orchestration RAG & LangChain & Consignes et appel des modèles \\
\hline
Interface web & Streamlit & Chatbot, recherche, fiches \\
\hline
Bibliothèques & Python & neo4j, langchain, streamlit, rdflib, etc. \\
\hline
\end{tabular}
\end{table}

% ------------------------------------------------------------
% II.1.5. Tableau de correspondance Matériel - Phases
% ------------------------------------------------------------

\subsection{Tableau de correspondance Matériel - Phases}
\label{subsec:correspondance_materiel}

Le tableau \ref{tab:correspondance_materiel} met en regard le matériel et les phases méthodologiques, et indique à quel moment et dans quel but chaque ressource entre en jeu.

\begin{table}[H]
\centering
\caption{Correspondance entre le matériel et les phases}
\label{tab:correspondance_materiel}
\begin{tabular}{|L{4cm}|L{3cm}|L{6cm}|}
\hline
\textbf{Matériel} & \textbf{Phase} & \textbf{Utilisation} \\
\hline
Corpus documentaire & Phase 3 & Constitution du corpus \\
\hline
Ontologies candidates & Phase 1 & Comparaison et sélection \\
\hline
gpt-6-luna & Phase 5 & Extraction \\
\hline
Llama 3 8B & Phase 7, Phase 8 & Génération RAG, évaluation \\
\hline
Google Colab & Phase 7, Phase 8 & Exécution de Llama 3 8B \\
\hline
Neo4j Desktop & Phase 6 & Base de connaissances \\
\hline
LangChain & Phase 7 & Orchestration RAG \\
\hline
Streamlit & Phase 7 & Interface web \\
\hline
Bibliothèques Python & Phase 5, 6, 7 & Développement \\
\hline
\end{tabular}
\end{table}

% ============================================================
% II.2. MÉTHODOLOGIE
% ============================================================

\section{Méthodologie}
\label{sec:methodologie}

Nous avons organisé la démarche en huit phases opérationnelles. Nous avons comparé les ontologies, défini les critères de recherche et constitué le corpus. Nous avons adapté l'ontologie, extrait les connaissances, construit la base de connaissances, développé le système d'accès et l'application web, puis évalué l'ensemble. Nous avons choisi l'ontologie avant de constituer le corpus afin de fixer le cadre conceptuel et de guider l'extraction.

% ------------------------------------------------------------
% II.2.1. Phase 1 : Identification et étude comparative des ontologies
% ------------------------------------------------------------

\subsection{Phase 1 : Identification et étude comparative des ontologies}
\label{subsec:phase1}

\textbf{Objectif} : Sélectionner l'ontologie la plus pertinente pour notre contexte.

Nous avons comparé les ontologies en trois temps. Nous avons défini une grille de critères, attribué une note à chacune, puis retenu celle qui a obtenu le meilleur score. Ce choix a fixé le modèle conceptuel des étapes suivantes.

\subsubsection{Définition de la grille de critères}

Nous avons retenu huit critères adaptés à nos besoins : couverture du domaine, richesse sémantique, disponibilité, documentation, interopérabilité, adaptabilité au contexte burkinabè, traçabilité et prise en compte des savoirs locaux. Le tableau \ref{tab:grille_criteres} présente ces critères.

\begin{table}[H]
\centering
\caption{Grille de critères de comparaison}
\label{tab:grille_criteres}
\begin{tabular}{|C{0.5cm}|L{3cm}|L{6.6cm}|C{2.5cm}|}
\hline
\textbf{N°} & \textbf{Critère} & \textbf{Description} & \textbf{Pondération} \\
\hline
1 & Couverture du domaine & Capacité à représenter plantes, usages, propriétés, maladies & Élevée (3) \\
\hline
2 & Richesse sémantique & Nombre et qualité des classes, relations, axiomes & Élevée (3) \\
\hline
3 & Disponibilité & Ontologie accessible et téléchargeable & Élevée (3) \\
\hline
4 & Documentation & Qualité de la documentation et des exemples & Moyenne (2) \\
\hline
5 & Interopérabilité & Compatibilité avec OWL, RDF, SPARQL & Moyenne (2) \\
\hline
6 & Adaptabilité & Facilité d'adaptation au contexte burkinabè & Élevée (3) \\
\hline
7 & Traçabilité & Capacité à associer les connaissances aux sources & Élevée (3) \\
\hline
8 & Prise en compte des savoirs locaux & Intégration des noms vernaculaires, recettes multi-plantes & Élevée (3) \\
\hline
\end{tabular}
\end{table}

Nous avons attribué à chaque critère une pondération proportionnelle à son importance. Les critères fortement pondérés (3) étaient déterminants pour représenter le contexte burkinabè, tandis que les critères moyennement pondérés (2) autorisaient davantage de compromis. Nous avons calculé le score global en multipliant chaque note par la pondération correspondante, en additionnant les produits et en divisant cette somme par la somme des pondérations. Le résultat est rapporté sur 5.

\subsubsection{Notation des ontologies}

Nous avons noté chaque ontologie de 1 (très faible) à 5 (très élevé) pour chacun des critères. Le tableau \ref{tab:notation_ontologies} présente les résultats.

\begin{table}[H]
\centering
\caption{Notation des ontologies candidates}
\label{tab:notation_ontologies}
\begin{tabular}{|L{3.5cm}|C{0.8cm}|C{0.8cm}|C{0.8cm}|C{0.8cm}|C{0.8cm}|C{0.8cm}|C{0.8cm}|C{0.8cm}|C{1.5cm}|}
\hline
\textbf{Ontologie} & \textbf{C1} & \textbf{C2} & \textbf{C3} & \textbf{C4} & \textbf{C5} & \textbf{C6} & \textbf{C7} & \textbf{C8} & \textbf{Score} \\
\hline
WATRIMed & 4 & 4 & 2 & 3 & 5 & 4 & 3 & 5 & 3.7 \\
\hline
ontoMEDTRAD & 5 & 5 & 5 & 4 & 4 & 5 & 4 & 5 & 4.7 \\
\hline
ATHMed & 4 & 4 & 3 & 3 & 4 & 3 & 2 & 3 & 3.2 \\
\hline
Deep Ontologies & 3 & 4 & 2 & 2 & 3 & 3 & 2 & 4 & 2.9 \\
\hline
FLOPO / PO & 2 & 3 & 5 & 4 & 5 & 2 & 1 & 1 & 2.7 \\
\hline
\end{tabular}
\end{table}

Nous avons analysé les notes afin de comparer la couverture du domaine et la prise en compte des savoirs locaux par les ontologies africaines. Nous avons également examiné la disponibilité et l'interopérabilité des ontologies botaniques, ainsi que les différences entre WATRIMed et ontoMEDTRAD concernant la richesse du modèle et l'accessibilité de la documentation.

Nous avons accordé une attention particulière à la disponibilité (C3) et à la documentation (C4). Nous avons évalué les ressources effectivement accessibles. WATRIMed est décrite par deux articles de congrès, dont l'un compte deux pages, mais son fichier d'ontologie n'est pas téléchargeable. La thèse consacrée à ontoMEDTRAD fournit la liste des classes et des propriétés, les métriques ainsi que l'origine des concepts.

\subsubsection{Sélection de l'ontologie}

WATRIMed \cite{watrimed2019, watrimed_design_2019} pourrait se distinguer par la richesse de son modèle documenté, qui compte 472 concepts et 75 propriétés, ainsi que par son alignement sur l'ontologie fondationnelle BioTopLite2. Toutefois, les deux publications n'en décrivent que les classes supérieures et les relations principales, et le fichier de l'ontologie n'est pas disponible pour une réutilisation directe.

Nous avons retenu ontoMEDTRAD \cite{kouame_systeme_2018}, qui a obtenu le meilleur score. Cette ontologie couvre l'Afrique de l'Ouest, représente les recettes multi-plantes, leurs ingrédients, les parties utilisées et les modes de préparation, et documente l'origine de ses concepts à partir de quinze ressources terminologiques et ontologiques. Sa thèse est accessible. Nous avons aussi pris en compte ses limites : sa validation globale reste annoncée comme perspective et son jeu de données comprend 22 plantes et 31 remèdes antipaludéens \cite{kouame_visual_2020}.

L'ontologie retenue fournira le cadre conceptuel pour extraire et organiser les connaissances. Elle guidera la constitution du corpus et son adaptation au contexte burkinabè.

\textbf{Livrable} : Ontologie ontoMEDTRAD sélectionnée.

% ------------------------------------------------------------
% II.2.2. Phase 2 : Définition des critères de recherche et de sélection
% ------------------------------------------------------------

\subsection{Phase 2 : Définition des critères de recherche et de sélection des documents}
\label{subsec:phase2}

\textbf{Objectif} : Définir les critères qui guideront la recherche et la sélection des documents scientifiques.

Nous avons défini deux catégories de critères, l'une pour rechercher les informations et l'autre pour sélectionner les documents. Elles visaient à assurer la pertinence, la qualité et la cohérence des sources retenues.

\subsubsection{Critères de recherche}

Nous avons retenu les critères suivants pour guider la recherche :

\begin{itemize}
\item \textbf{Les plantes} : espèces végétales documentées au Burkina Faso, avec identification botanique (nom scientifique, famille, synonymes).
\item \textbf{Les parties utilisées} : feuilles, racines, écorces, fruits, graines, fleurs, tiges, etc.
\item \textbf{Les pathologies et usages médicinaux} : paludisme, infections respiratoires, digestives, diabète, hypertension, etc.
\item \textbf{Les modes de préparation} : décoction, macération, infusion, poudre, extrait alcoolique, etc.
\item \textbf{Les niveaux de preuve} : savoirs traditionnels, études in vitro, études in vivo, essais cliniques.
\item \textbf{Les propriétés pharmacologiques} : antimicrobiennes, antipaludiques, anti-inflammatoires, antioxydantes, toxicité, etc.
\item \textbf{La composition chimique} : alcaloïdes, flavonoïdes, tanins, saponines, huiles essentielles, etc.
\end{itemize}

Ils couvrent les principales dimensions des connaissances sur les plantes médicinales et ont permis de cibler les informations utiles à l'étude.

\subsubsection{Critères de sélection des documents}

Nous avons sélectionné les documents selon les critères suivants :

\begin{itemize}
\item \textbf{Pertinence} : le document traite des plantes médicinales du Burkina Faso.
\item \textbf{Identification de la plante} : le nom scientifique est précisé.
\item \textbf{Disponibilité d'une référence scientifique} : le document est publié dans une revue à comité de lecture, une thèse, un mémoire ou un rapport de recherche.
\item \textbf{Fiabilité} : les sources sont vérifiables et citées.
\item \textbf{Actualité} : les documents récents sont privilégiés, sans exclure les travaux fondateurs.
\end{itemize}

Ces critères nous ont permis d'écarter les documents non scientifiques ou non pertinents et de constituer un corpus cohérent pour l'étude.

\textbf{Livrable} : grille de critères de recherche et de sélection.

% ------------------------------------------------------------
% II.2.3. Phase 3 : Constitution du corpus scientifique
% ------------------------------------------------------------

\subsection{Phase 3 : Constitution du corpus scientifique}
\label{subsec:phase3}

\textbf{Objectif} : Rassembler un corpus documentaire fiable et pertinent sur les plantes médicinales du Burkina Faso.

Nous avons constitué le corpus en quatre étapes. Nous avons recherché les documents dans Google Scholar, PubMed et ResearchGate, puis consulté burkinaherbiermed.org et les données de l'INERA. Nous avons inclus les documents portant sur les plantes médicinales du Burkina Faso, précisant le nom scientifique des plantes et publiés ou validés par une institution scientifique avec des sources vérifiables. Nous avons écarté les documents non scientifiques, ceux sans identification botanique et les documents anciens sans valeur de référence. Nous avons classé les documents par type, plante, pathologie et région et les avons enregistrés dans un répertoire structuré. Le corpus comprend 35 documents scientifiques et les données de burkinaherbiermed.org et de l'INERA.

\textbf{Livrable} : un corpus de 35 documents scientifiques, complété par les données de burkinaherbiermed.org et de l'INERA.

% ------------------------------------------------------------
% II.2.4. Phase 4 : Adaptation de l'ontologie sélectionnée
% ------------------------------------------------------------

\subsection{Phase 4 : Adaptation de l'ontologie sélectionnée}
\label{subsec:phase4}

\textbf{Objectif} : Adapter le modèle conceptuel de l'ontologie sélectionnée au contexte burkinabè.

Nous avons examiné les concepts existants, ajouté des classes et des propriétés, modifié certaines relations et documenté les changements.

\subsubsection{Analyse des concepts existants}

Nous avons étudié la structure de l'ontologie ontoMEDTRAD telle qu'elle est décrite dans la thèse \cite{kouame_systeme_2018}. L'ontologie compte 144 concepts, 76 propriétés associatives, 16 propriétés de données, 789 individus et 2 734 axiomes. Elle est formalisée en logique de description ALCHOIQ, éditée sous Protégé~5 et vérifiée par les raisonneurs HermiT et Pellet.

Le fichier OWL de l'ontologie n'étant pas téléchargeable publiquement, nous avons reconstitué sa structure à partir de la thèse. Nous nous sommes appuyés sur les captures d'écran de Protégé (figures 5.15 à 5.20) et les tableaux récapitulatifs des concepts, propriétés et métriques (tableaux 5.5, 5.21 et 5.22). Cette reconstitution a servi de base à l'adaptation présentée au chapitre \ref{chap:conception_developpement}.

Nous avons reconstitué la structure conceptuelle de l'ontologie, ses classes, propriétés et relations, mais pas ses instances d'origine. Nous n'avons pas réutilisé les 789 individus d'ontoMEDTRAD, liés au contexte et aux sources des travaux d'origine. Cette phase a produit le modèle conceptuel adapté ; le peuplement par les individus extraits de la littérature scientifique est intervenu en phase 5.

Ses concepts principaux décrivent le parcours de soin, du praticien et du patient aux symptômes, aux maladies, aux ressources médicinales, à l'examen clinique, aux recettes, aux ingrédients, aux parties utilisées, aux préparations et aux remèdes. L'ontologie distingue cinq types de plantes et neuf parties utilisées. Elle comprend aussi une classe pour les références scientifiques associées aux ressources, aux recettes et aux remèdes. Les concepts secondaires couvrent les contre-indications, les propriétés thérapeutiques, la toxicité, l'étiologie et la classification des maladies selon la CIM-10. Les relations de base expriment la subsomption et la relation partie-tout. Soixante-seize propriétés associatives complètent le modèle, dont la hiérarchie \textit{aPourForme}, utilisée pour l'iconisation des plantes.

Nous avons ensuite évalué ce modèle au regard du contexte burkinabè. Il couvre les concepts de base et les raisonneurs testés n'y ont signalé aucune incohérence, mais il nécessitait des adaptations pour représenter les noms vernaculaires en langues locales, les maladies spécifiques, les niveaux de preuve, les dosages et la traçabilité détaillée des sources. Sa validation globale reste une perspective, et son jeu de données se limite à 22 plantes et 31 remèdes antipaludéens.

\subsubsection{Ajout de classes}

La comparaison des fichiers du modèle conceptuel adapté et du schéma conservé dénombre 40 classes ajoutées. Le modèle adapté déclare également 25 propriétés objet au total, dont 20 absentes du schéma conservé, et 20 propriétés de données, dont 15 ajoutées. Les classes ajoutées couvrent notamment les noms vernaculaires, les usages thérapeutiques, les niveaux de preuve, les sources, les préparations, des parties végétales complémentaires, les composés chimiques, les activités pharmacologiques et les effets toxicologiques. Ces ajouts définissent le modèle conceptuel ; les individus et données extraits sont intégrés lors de la phase suivante.

Ces ajouts structurent des connaissances locales et scientifiques à différents niveaux. Ils ne préchargent pas les valeurs, comme les noms, pathologies, régions ou composés ; celles-ci proviennent des sources du corpus et sont enregistrées dans la base avec leur provenance.

\subsubsection{Ajout de propriétés}

Nous avons ajouté des propriétés de données pour les noms locaux, le niveau de preuve, le type d'étude, les références, le dosage et la toxicité, en complément de propriétés textuelles pour les noms vernaculaires et les descriptions d'activités ou d'effets. Des propriétés objet relient également l'usage à sa maladie cible, à la partie végétale utilisée, à la préparation et à la source ; d'autres décrivent les associations avec des composés, activités pharmacologiques, risques et localités. Les noms techniques exacts sont ceux déclarés dans la version courante du fichier Turtle.

L'enrichissement ainsi obtenu saisit des informations fines (noms vernaculaires, coordonnées, niveaux de preuve, références précises) que l'ontologie d'origine ne documentait pas pour le Burkina Faso.

\subsubsection{Modification de relations}

Nous avons structuré les liens entre plante, usage, maladie, partie utilisée et préparation, puis généralisé la traçabilité documentaire. La propriété objet \textit{estDocumenteePar} relie les ressources du graphe à une source scientifique; des sous-propriétés spécialisées permettent de préciser si la source documente une plante, un usage, un nom vernaculaire, une activité, un composé ou un effet toxicologique. \textit{aPourNiveauDePreuve} rattache un usage au niveau déclaré par la source. Ces liens ne signifient pas que les usages ou activités constituent des recommandations thérapeutiques ni une validation clinique.

Ces modifications assurent la traçabilité des connaissances et permettent d'exprimer des relations que le modèle d'origine, fondé sur la relation « est un » et sur la relation partie-tout, ne couvrait pas.

\subsubsection{Documentation des adaptations}

Nous avons consigné toutes les modifications dans un journal : chaque ajout est justifié, chaque choix est explicité, et les classes, propriétés et relations ajoutées, modifiées ou supprimées y sont décrites une à une.

Cette documentation n'est pas une formalité. Elle rend la démarche reproductible par un tiers et facilite les évolutions ultérieures de l'ontologie.

\textbf{Livrable} : Le modèle conceptuel adapté de l'ontologie, sans les données instanciées.

% ------------------------------------------------------------
% II.2.5. Phase 5 : Enrichissement de l'ontologie
% ------------------------------------------------------------

\subsection{Phase 5 : Enrichissement des connaissances à partir de la littérature scientifique}
\label{subsec:phase5}

\textbf{Objectif} : Extraire les connaissances du corpus pour alimenter la base organisée selon l'ontologie.

À partir du modèle conceptuel adapté en phase 4, nous avons enrichi l'ontologie avec des individus et des données extraits des sources scientifiques. Nous avons préparé ces données structurées et sourcées pour leur exploitation dans la base orientée graphe Neo4j.

Nous avons conçu les consignes d'extraction, procédé à l'extraction assistée par modèle de langue, contrôlé les données obtenues, puis structuré et préparé ces données pour leur intégration dans la base.

\subsubsection{Conception des prompts d'extraction}

Nos consignes d'extraction ciblaient les noms scientifiques et vernaculaires, la famille botanique, les parties utilisées, les usages, les préparations, les activités pharmacologiques, le niveau de preuve et la source. Le chapitre \ref{chap:conception_developpement} présente un exemple de consigne, avec le format JSON et les champs demandés.

Nous avons employé une même consigne pour les documents traités. Les champs absents reçoivent la valeur \texttt{null} ou une liste vide, et aucune valeur n'est complétée par déduction.

\subsubsection{Exécution de l'extraction}

Nous avons soumis les documents à gpt-6-luna dans l'environnement OpenWork, un par un, et enregistré les résultats en JSON. L'extraction a porté sur dix des 35 documents du corpus, principalement des enquêtes ethnobotaniques ; les 25 autres n'ont pas été extraits.

\subsubsection{Contrôle de la qualité des extractions}

Nous avons repris le niveau de preuve indiqué par la source lorsqu'il était explicite, qu'il s'agisse d'un savoir traditionnel, d'une étude \textit{in vitro}, d'une étude \textit{in vivo} ou d'un essai clinique. Nous n'avons pas évalué indépendamment la solidité des faits et avons laissé le champ vide lorsque la source ne précisait pas ce niveau. Un expert externe n'a pas relu les extractions. Nous avons contrôlé leur structure, puis comparé une partie des faits aux documents d'origine à la phase 8.

\subsubsection{Structuration des données}

Nous avons organisé les données selon le modèle de l'ontologie dans un classeur d'individus, par type de connaissance (sources, plantes, noms vernaculaires, usages, préparations, maladies, activités, composés, toxicologie et localités), puis les avons transformées en individus au format Turtle.

\subsubsection{Préparation de l'intégration dans la base de connaissances}

Le fichier OntoPlantMedBF-peuplee.ttl contient des individus extraits de dix documents scientifiques, en plus du schéma adapté. Ce fichier constitue l'ontologie enrichie importée dans Neo4j à la phase 6. Les données de burkinaherbiermed.org et de l'INERA, consultées lors de la constitution du corpus, n'y ont pas été intégrées.

\textbf{Livrable} : OntoPlantMedBF, ontologie enrichie avec les données extraites et leurs sources.

% ------------------------------------------------------------
% II.2.6. Phase 6 : Construction et interrogation de la base
% ------------------------------------------------------------

\subsection{Phase 6 : Construction et interrogation de la base de connaissances}
\label{subsec:phase6}

\textbf{Objectif} : Importer et exploiter dans Neo4j l'ontologie réajustée \textit{OntoPlantMedBF}.

Nous avons importé dans Neo4j l'ontologie réajustée \textit{OntoPlantMedBF} et ses données, en les représentant sous forme de nœuds, propriétés et relations du graphe. Nous avons vérifié la correspondance avec le modèle ontologique et testé les requêtes de la base.

\subsubsection{Import et ajustement du modèle dans le graphe}

Nous avons importé les classes et propriétés d'OntoPlantMedBF dans Neo4j sous forme de labels, propriétés et relations. Les principaux nœuds représentent les plantes, maladies, usages, parties, préparations, sources, composés chimiques et régions. Les relations du graphe reprennent les liens du modèle, avec des noms adaptés aux conventions Neo4j, et les propriétés conservent notamment les noms, niveaux de preuve, DOI et coordonnées géographiques.

Nous avons ainsi importé dans le graphe la structure réajustée d'OntoPlantMedBF et ses données. Les entités sont représentées par des nœuds, leurs liens par des relations et leurs caractéristiques par des propriétés, conformément au modèle ontologique.

\subsubsection{Import des données}

Nous avons chargé les données préparées, créé ou rapproché les nœuds et relations au moyen de requêtes Cypher, puis vérifié l'intégrité et la cohérence de l'import avec le modèle ontologique.

Après l'import, nous avons contrôlé que les données, les relations et les propriétés avaient été correctement reprises.

\subsubsection{Requêtes de test}

Nous avons testé quatre familles de requêtes : les requêtes simples sur les plantes associées à une maladie, les requêtes complexes sur les recettes multi-plantes, les requêtes de traçabilité et les requêtes spatiales par région.

Ces vérifications nous ont permis de contrôler le fonctionnement de la base et les résultats renvoyés par les requêtes.

\textbf{Livrable} : Une base Neo4j peuplée et interrogeable, complémentaire d'OntoPlantMedBF.

% ------------------------------------------------------------
% II.2.7. Phase 7 : Exploitation par l'IA et développement web
% ------------------------------------------------------------

\subsection{Phase 7 : Exploitation par IA et développement de l'application web}
\label{subsec:phase7}

\textbf{Objectif} : Développer le système d'accès intelligent et l'interface utilisateur.

Nous avons mis en place le système RAG en définissant son architecture, en construisant la chaîne de traitement et en intégrant ses composants.

\subsubsection{Architecture du système}

L'architecture relie l'utilisateur à Streamlit, LangChain, Neo4j et Llama 3 jusqu'à la réponse. Elle associe la recherche dans la base à la génération par le modèle de langue et conserve le lien entre les informations et leurs sources. Le chapitre \ref{chap:conception_developpement} décrit la chaîne de traitement, de l'interprétation de la question à l'affichage des références.

\subsubsection{Intégration des composants}

Nous avons connecté l'application à Neo4j et au modèle de langue, puis configuré les consignes de génération des réponses. Cette intégration a réuni les composants du système RAG.

\subsubsection{Développement de l'application web}

Nous avons conçu l'application autour d'une page d'accueil avec chatbot, d'une recherche par plante, d'une recherche par pathologie et de fiches descriptives. Nous y avons intégré le système RAG et la gestion de l'historique. L'interface permet d'interroger la base en langage naturel et de recevoir des réponses reliées à leurs sources. Le chapitre \ref{chap:conception_developpement} en décrit la conception et les composants.

\subsubsection{Tests et validation}

Nous avons mené des tests fonctionnels pour vérifier le fonctionnement de l'application et des tests utilisateurs pour examiner sa facilité d'utilisation et la pertinence des réponses.

\textbf{Livrable} : Un système RAG fonctionnel et une application web fonctionnelle.

% ------------------------------------------------------------
% II.2.8. Phase 8 : Évaluation et validation
% ------------------------------------------------------------

\subsection{Phase 8 : Évaluation et validation}
\label{subsec:phase8}

\textbf{Objectif} : Évaluer la qualité de l'ontologie, du système RAG et de l'application web.

Nous avons évalué trois composantes : l'ontologie, le système RAG et l'application web.

\subsubsection{Évaluation de l'ontologie}

Nous avons évalué l'ontologie selon les quatre critères présentés dans le tableau \ref{tab:eval_ontologie}.

\begin{table}[H]
\centering
\caption{Critères d'évaluation de l'ontologie}
\label{tab:eval_ontologie}
\begin{tabular}{|L{3cm}|L{6cm}|L{4cm}|}
\hline
\textbf{Critère} & \textbf{Description} & \textbf{Méthode} \\
\hline
Couverture & Capacité à représenter les connaissances du corpus & Vérification manuelle \\
\hline
Cohérence & Absence de contradictions & Requêtes de test \\
\hline
Traçabilité & Association des connaissances aux sources & Vérification des relations \\
\hline
Extensibilité & Facilité d'ajout de nouvelles connaissances & Test d'ajout \\
\hline
\end{tabular}
\end{table}

\subsubsection{Évaluation du système RAG}

Nous avons évalué le système RAG selon les quatre critères présentés dans le tableau \ref{tab:eval_rag}.

\begin{table}[H]
\centering
\caption{Critères d'évaluation du système RAG}
\label{tab:eval_rag}
\begin{tabular}{|L{3cm}|L{6cm}|L{4cm}|}
\hline
\textbf{Critère} & \textbf{Description} & \textbf{Métrique} \\
\hline
Pertinence & Les réponses répondent à la question & Précision \\
\hline
Fidélité & Les réponses sont fidèles aux sources & Rappel \\
\hline
Exactitude & Les réponses sont correctes & F1-score \\
\hline
Traçabilité & Les sources sont affichées & Taux de couverture \\
\hline
\end{tabular}
\end{table}

Nous avons constitué un jeu de 50 questions de référence portant sur les plantes, les usages et leurs sources, puis comparé les réponses du système aux réponses attendues.

\subsubsection{Évaluation de l'application web}

Nous avons évalué l'application web selon les quatre critères présentés dans le tableau \ref{tab:eval_web}.

\begin{table}[H]
\centering
\caption{Critères d'évaluation de l'application web}
\label{tab:eval_web}
\begin{tabular}{|L{3cm}|L{6cm}|L{4cm}|}
\hline
\textbf{Critère} & \textbf{Description} & \textbf{Méthode} \\
\hline
Facilité d'utilisation & L'interface est intuitive & Questionnaire \\
\hline
Compréhension & Les réponses sont claires & Questionnaire \\
\hline
Utilité & L'application répond aux besoins & Questionnaire \\
\hline
Performance & Temps de réponse & Mesure \\
\hline
\end{tabular}
\end{table}

Nous avons sollicité des étudiants, des chercheurs et des tradipraticiens pour cette évaluation.

\textbf{Livrable} : Rapport d'évaluation et de validation.

% ============================================================
% II.3. SYNTHÈSE ET TRANSITION
% ============================================================

\section{Synthèse et transition}
\label{sec:synthese}

\subsection{Tableau de correspondance Phases - Livrables}

Le tableau \ref{tab:correspondance_phases} récapitule les objectifs et les livrables des huit phases.

\begin{table}[H]
\centering
\caption{Correspondance entre les phases, les objectifs et les livrables}
\label{tab:correspondance_phases}
\begin{tabular}{|L{2.2cm}|L{4.5cm}|L{6.8cm}|}
\hline
\textbf{Phase} & \textbf{Objectif} & \textbf{Livrable} \\
\hline
Phase 1 & Sélectionner l'ontologie la plus pertinente & Ontologie ontoMEDTRAD sélectionnée \\
\hline
Phase 2 & Définir les critères de recherche et de sélection & Grille de critères \\
\hline
Phase 3 & Rassembler un corpus fiable et pertinent & Corpus de 35 documents \\
\hline
Phase 4 & Adapter le modèle conceptuel de l'ontologie & Modèle conceptuel adapté \\
\hline
Phase 5 & Extraire les connaissances et enrichir l'ontologie & Ontologie enrichie en données \\
\hline
Phase 6 & Importer OntoPlantMedBF dans la base graphe & OntoPlantMedBF intégrée à Neo4j, base interrogeable \\
\hline
Phase 7 & Développer le système RAG et l'application web & Système RAG et application web \\
\hline
Phase 8 & Évaluer la qualité du système & Rapport d'évaluation \\
\hline
\end{tabular}
\end{table}

% ============================================================
% CONCLUSION DU CHAPITRE
% ============================================================

\section*{Conclusion du chapitre}
\addcontentsline{toc}{section}{Conclusion du chapitre}

Ce chapitre a exposé le matériel et la méthodologie de notre recherche. Le matériel réunit un corpus de 35 documents scientifiques, les données de burkinaherbiermed.org et de l'INERA, cinq ontologies candidates et un environnement technique bâti sur gpt-6-luna pour l'extraction, Llama 3 8B sur Google Colab pour le RAG, Neo4j Desktop, LangChain et Streamlit.

Nous avons suivi huit phases : comparaison des ontologies, définition des critères de recherche et de sélection, constitution du corpus, adaptation de l'ontologie, extraction des connaissances, construction et interrogation de la base, développement du système d'accès et de l'application, puis évaluation. Nous avons décrit les actions menées et les livrables associés.

Nous avons choisi l'ontologie avant de constituer le corpus afin d'assurer la cohérence entre le modèle conceptuel et les données collectées, et de cibler les informations pertinentes dès la collecte.

Nous signalons deux limites. Les niveaux de preuve associés aux connaissances proviennent des sources et n'ont pas fait l'objet d'une évaluation indépendante. De plus, un expert externe n'a pas validé les extractions. Ces limites sont à prendre en compte dans les travaux ultérieurs. Nous avons néanmoins conservé le lien entre chaque information et sa source.

Le chapitre \ref{chap:resultats_discussion} présente les résultats de l'évaluation menée en phase 8.

Le chapitre suivant présente la conception et le développement du système, notamment l'ontologie adaptée, la base de connaissances, le système RAG et l'application web.
).

\phantomsection
\section*{Introduction du chapitre}

Après avoir posé le cadre conceptuel et exposé l'état de l'art, nous présentons le matériel mobilisé et la démarche méthodologique suivie. Ce chapitre fait le lien entre la réflexion théorique et la mise en œuvre technique.

Nous présentons le matériel mobilisé, notamment le corpus documentaire, les ontologies candidates et l'environnement technique. Nous décrivons les huit phases de la démarche suivie, de l'étude comparative des ontologies à l'évaluation du système, puis en proposons une synthèse.

% ============================================================
% II.1. MATÉRIEL
% ============================================================

\section{Matériel}
\label{sec:materiel}

Le matériel comprend le corpus documentaire, les ontologies candidates et l'environnement technique. Les deux tableaux qui suivent récapitulent ces ressources, réparties entre documentation, modèles conceptuels et outils techniques.

% ------------------------------------------------------------
% II.1.1. Corpus documentaire
% ------------------------------------------------------------

\subsection{Corpus documentaire}
\label{subsec:corpus}

Nous avons choisi l'ontologie avant de constituer le corpus afin de déterminer les notions et les relations qui guideraient l'extraction. Nous avons réuni les publications scientifiques et les données institutionnelles consultées sur burkinaherbiermed.org et auprès de l'INERA.

La figure \ref{fig:flux_corpus} résume le processus de constitution du corpus, du choix de l'ontologie à la sélection des documents.

\begin{figure}[H]
\centering
\includegraphics[width=0.8\textwidth]{flux_corpus.png}
\caption{Constitution du corpus documentaire}
\label{fig:flux_corpus}
\end{figure}

\subsubsection{Sources scientifiques publiées}

Notre corpus scientifique comprend 35 rapports de recherche, articles, thèses et mémoires provenant de revues, d'universités et d'institutions de recherche burkinabè. Ces travaux ethnobotaniques, ethnopharmacologiques et phytochimiques documentent les noms scientifiques et vernaculaires, les parties utilisées, les pathologies traitées, les modes de préparation et, pour les études qui les rapportent, les dosages et les données de toxicité.

Nous avons retenu ces travaux pour leur qualité et leur pertinence au contexte burkinabè, des publications fondatrices aux plus récentes. Cette diversité d'espèces, de pathologies et de pratiques thérapeutiques soutient l'extraction. Quatre travaux illustrent les types de données mobilisés dans le corpus :
\begin{itemize}
\item Zerbo et al.~\cite{zerbo_plantes_2011} ont étudié les plantes médicinales chez les Sanan du Burkina Faso. Leur travail documente 85 espèces utilisées pour traiter 63 affections.
\item Ouédraogo et al.~\cite{ouedraogo_revue_2024} ont dressé une revue des plantes médicinales à potentiel aromatique du pays, recensant 42 espèces et en caractérisant les propriétés et les usages.
\item Tibiri et al.~\cite{tibiri2020} ont mené une enquête à l'échelle nationale sur les plantes utilisées contre les maladies hépatiques, documentant plus de 2 000 entrées de plantes correspondant à 212 espèces.
\item Ouédraogo et al.~\cite{ouedraogo2020markets} ont enquêté auprès de 30 commerçants de plantes médicinales dans six grandes villes, identifiant 104 espèces utilisées pour diverses pathologies.
\end{itemize}

Ces travaux fournissent des méthodes réutilisables et des observations de terrain recueillies au Burkina Faso.

\subsubsection{Données de l'Herbier National des Plantes Médicinales du Burkina Faso}

Le site burkinaherbiermed.org diffuse des données sur les plantes médicinales du pays, notamment des informations botaniques, des usages thérapeutiques et des références bibliographiques. Ces données proviennent du projet « Herbier avec chimiothèque », piloté par le Ministère de la Santé. Ce projet a pour objectif de constituer un herbier national dédié aux plantes médicinales et une chimiothèque de marqueurs chimiques. Il documente les plantes utilisées contre les dix pathologies prioritaires définies par le ministère, géolocalise les stations et conserve des spécimens. Ces informations alimentent une base nationale accessible par application mobile et site web. Le projet associe un herbier, une chimiothèque et une base de données ethnobotaniques et géographiques. Il complète les autres herbiers burkinabè, notamment le HNBU du CNRST/INERA, le Centre National de Semences Forestières et l'herbier de l'Université Joseph Ki-Zerbo.

\subsubsection{Données de l'INERA}

La base de données de l'Institut de l'Environnement et de Recherches Agricoles (INERA) comprend plus de 7 000 espèces végétales et les coordonnées géographiques des stations où elles ont été observées. Ces données complètent les informations ethnobotaniques par une dimension spatiale et permettent d'étudier la répartition des plantes au Burkina Faso. Le rapprochement des usages rapportés et de la distribution géographique peut faire apparaître des régularités régionales, qui restent à vérifier.

Le tableau \ref{tab:corpus} récapitule les sources du corpus.

\begin{table}[H]
\centering
\caption{Sources du corpus documentaire}
\label{tab:corpus}
\begin{tabular}{|L{4cm}|L{2.8cm}|L{3cm}|L{5cm}|}
\hline
\textbf{Source} & \textbf{Type} & \textbf{Volume estimé} & \textbf{Contenu principal} \\
\hline
Articles et thèses & Publications scientifiques & 35 documents & Études ethnobotaniques, pharmacologiques \\
\hline
burkinaherbiermed.org & Base de données en ligne & Plusieurs centaines de plantes & Données botaniques, usages, références \\
\hline
INERA & Base de données institutionnelle & > 7 000 espèces & Coordonnées géographiques, identification \\
\hline
\end{tabular}
\end{table}

Les documents du corpus relèvent de types d'étude différents, et chacun est exploitable pour les informations qu'il rapporte. Trois profils se distinguent. Les enquêtes ethnobotaniques et ethnopharmacologiques, menées auprès de tradipraticiens ou dans les marchés, rapportent les noms vernaculaires, les parties utilisées, les modes de préparation et les pathologies traitées, souvent plante par plante. Les études expérimentales, notamment les essais \textit{in vitro} et les travaux phytochimiques, mesurent des activités pharmacologiques, des concentrations, des niveaux de preuve et parfois les lieux de récolte. Elles renseignent moins sur les noms vernaculaires et les pratiques de préparation. Enfin, les revues de synthèse dressent des inventaires à grande échelle et servent surtout à repérer les sources primaires.

Nous avons appliqué trois règles à l'extraction. Chaque source ne renseignait que les champs qu'elle mesurait. Nous avons laissé vides les champs non couverts, sans les compléter par déduction. Nous n'avons pas retenu comme fait établi un usage traditionnel rapporté dans un article expérimental sans référence ethnobotanique. Enfin, nous n'avons pas additionné les nombres d'espèces d'une source à l'autre, car une revue peut reprendre les espèces documentées dans les enquêtes qu'elle cite. Le tableau \ref{tab:corpus} présente les sources document par document, sans agréger les volumes.

Ce corpus rassemble des études menées dans plusieurs régions du Burkina Faso, portant sur différentes pathologies et mobilisant des approches méthodologiques variées. Il a servi de base à l'extraction des connaissances et à l'enrichissement de l'ontologie. Le choix préalable de l'ontologie a guidé l'extraction et l'organisation des documents selon les concepts et les relations définis.

% ------------------------------------------------------------
% II.1.2. Ontologies candidates
% ------------------------------------------------------------

\subsection{Ontologies candidates}
\label{subsec:ontologies_candidates}

Nous avons retenu cinq ontologies candidates pour l'étude comparative. Le chapitre \ref{chap:cadre_theorique}, aux sections \ref{subsubsec:ontologies_internationales} et \ref{subsubsec:ontologies_africaines}, les présente. Elles couvrent différents domaines, notamment la botanique, la médecine traditionnelle, la pharmacologie et l'ethnomédecine. Le tableau \ref{tab:ontologies_candidates} en présente les caractéristiques.

Nous les avons comparées à l'aide de la grille de critères présentée en phase 1 (section \ref{subsec:phase1}). Nous avons sélectionné l'ontologie la plus adaptée à notre contexte avant de constituer le corpus, afin de fixer le modèle conceptuel qui a guidé l'extraction.

\begin{table}[H]
\centering
\caption{Ontologies candidates pour la comparaison}
\label{tab:ontologies_candidates}
\begin{tabular}{|C{0.5cm}|L{3cm}|L{3cm}|L{3cm}|L{2cm}|L{2.5cm}|}
\hline
\textbf{N°} & \textbf{Ontologie} & \textbf{Domaine} & \textbf{Zone géographique} & \textbf{Format} & \textbf{Source d'accès} \\
\hline
1 & WATRIMed & Médecine traditionnelle & Afrique de l'Ouest & OWL, graphe de connaissances & IOS Press (2 articles) \\
\hline
2 & ontoMEDTRAD & Médecine traditionnelle & Afrique de l'Ouest & OWL + iconique & Thèse A. Kouamé (UGB) \\
\hline
3 & ATHMed & Médecine traditionnelle & Afrique (Ghana) & Cadre conceptuel & Frontiers in AI \\
\hline
4 & Deep Ontologies & Médecine traditionnelle & Afrique & OWL (Protégé) & Articles ARIMA \\
\hline
5 & FLOPO, PO & Botanique & International & OWL & OBO Foundry \\
\hline
\end{tabular}
\end{table}

Chacune présente des acquis et des lacunes. Les ontologies botaniques (FLOPO, PO) décrivent bien la plante mais ne traitent pas des usages médicinaux. Les ontologies pharmaceutiques (DrOn) se resserrent sur les médicaments modernes. L'OTM, prometteuse, reste sans ancrage africain. Les initiatives africaines (WATRIMed, ontoMEDTRAD, ATHMed) s'intéressent en revanche à la médecine traditionnelle du terrain, avec ses dimensions culturelles et ses recettes multi-plantes, même si ATHMed en reste, elle, au stade d'un cadre proposé. La phase 1 en livre la comparaison détaillée.

% ------------------------------------------------------------
% II.1.3. Outils et environnement technique
% ------------------------------------------------------------

\subsection{Outils et environnement technique}
\label{subsec:environnement}

Les outils techniques occupent une place importante dans notre démarche. Nous décrivons les phases concernées et les raisons de leur choix.

La figure \ref{fig:dispositif_technique} résume l'environnement d'expérimentation, en précisant où chaque composant est exécuté et comment ils communiquent.

\begin{figure}[H]
\centering
\includegraphics[width=0.8\textwidth]{dispositif_technique.png}
\caption{Dispositif technique d'expérimentation}
\label{fig:dispositif_technique}
\end{figure}

\subsubsection{Modèles de langage}

Nous avons utilisé deux modèles, chacun pour une étape distincte. Pour l'extraction des connaissances (phase 5), nous avons utilisé \textbf{gpt-6-luna} d'OpenAI, accessible par l'environnement OpenWork. Pour la génération des réponses du système RAG (phase 7), nous avons utilisé \textbf{Llama 3 8B} \cite{meta2024llama3}, un modèle ouvert. Gemini (\textit{gemini-3.5-flash-lite}) nous a servi de point de comparaison lors de l'évaluation.

Le tableau \ref{tab:llama_acces} présente les modes d'accès à ces modèles.

\begin{table}[H]
\centering
\caption{Modèles de langage et modes d'accès}
\label{tab:llama_acces}
\begin{tabular}{|L{4cm}|L{4cm}|L{6cm}|}
\hline
\textbf{Usage} & \textbf{Modèle} & \textbf{Mode d'accès} \\
\hline
Extraction (phase 5) & gpt-6-luna & OpenWork \\
\hline
Génération RAG (phase 7) & Llama 3 8B & Ollama sur Google Colab, API exposée par ngrok \\
\hline
Comparaison (phase 8) & gemini-3.5-flash-lite & API de Google \\
\hline
\end{tabular}
\end{table}

Nous avons retenu Llama 3 pour la génération parce qu'il s'agit d'un modèle ouvert que nous pouvons exécuter sur une instance que nous contrôlons. La version 8B tient dans la mémoire d'un GPU accessible gratuitement sur Colab. Ollama sert le modèle et ngrok expose son API à l'application web, qui reste sur le poste local.

\subsubsection{Plateforme d'exécution : Google Colab}

Nous avons exécuté Llama 3 8B sur \textbf{Google Colab}, une plateforme de notebooks donnant accès à des GPU adaptés à l'inférence.

\textbf{Phases d'utilisation :} phases 7 et 8 pour la génération des réponses et l'évaluation du système RAG.

Cette plateforme nous a donné accès à un GPU sans nécessiter l'acquisition d'une infrastructure locale. La session reste cependant temporaire et l'adresse exposée par ngrok change à chaque lancement, ce qui convient à un prototype mais pas à un service pérenne.

\subsubsection{Base de connaissances : Neo4j Desktop}

Nous avons utilisé \textbf{Neo4j Desktop}, la version locale de Neo4j, pour créer, gérer et interroger la base de connaissances orientée graphe.

\textbf{Phase d'utilisation :} phase 6 pour la construction et l'interrogation de la base.

Le tableau \ref{tab:neo4j} présente les caractéristiques de Neo4j Desktop.

\begin{table}[H]
\centering
\caption{Caractéristiques de Neo4j Desktop}
\label{tab:neo4j}
\begin{tabular}{|L{5cm}|L{8cm}|}
\hline
\textbf{Caractéristique} & \textbf{Valeur} \\
\hline
Type & Base de données orientée graphe \\
\hline
Langage de requête & Cypher \\
\hline
Version & Desktop (locale) \\
\hline
Format d'import & CSV, JSON \\
\hline
Interface & Neo4j Browser, Bloom \\
\hline
\end{tabular}
\end{table}

Nous avons choisi Neo4j pour représenter les relations entre les plantes, les usages, les propriétés et les sources, interroger le graphe avec Cypher et assurer la traçabilité des connaissances. La base permet également d'ajouter des connaissances et de visualiser les liens entre les entités.

\subsubsection{Orchestration RAG : LangChain}

Nous avons utilisé \textbf{LangChain} pour définir les consignes et appeler les modèles de langage de façon uniforme. La recherche dans Neo4j passe par le pilote officiel et par nos propres requêtes.

\textbf{Phase d'utilisation :} phase 7 pour l'accès aux connaissances et le développement web.

Le tableau \ref{tab:langchain} présente les composants LangChain utilisés.

\begin{table}[H]
\centering
\caption{Composants LangChain utilisés}
\label{tab:langchain}
\begin{tabular}{|L{5cm}|L{8cm}|}
\hline
\textbf{Composant} & \textbf{Rôle} \\
\hline
PromptTemplate & Préparation des consignes de génération \\
\hline
ChatOpenAI & Appel de Llama 3 par l'API d'Ollama, compatible avec celle d'OpenAI \\
\hline
ChatGoogleGenerativeAI & Appel de Gemini pour la comparaison \\
\hline
\end{tabular}
\end{table}

Nous avons retenu LangChain parce qu'il permet de changer de modèle sans modifier le reste du système. Nous n'avons pas confié la sélection des données à une chaîne générant du Cypher : le chapitre \ref{chap:conception_developpement} justifie ce choix \cite{lewis2020rag}.

\subsubsection{Interface web : Streamlit}

Nous avons développé l'application web avec \textbf{Streamlit}, un framework Python qui permet de créer des applications web interactives.

\textbf{Phase d'utilisation :} phase 7 pour le développement de l'application web.

L'application propose un chatbot en langage naturel, des recherches par plante, pathologie ou région, des fiches descriptives et l'affichage des sources et des niveaux de preuve. Nous avons également prévu une fonction de visualisation du graphe de connaissances.

Nous avons retenu Streamlit pour sa simplicité d'utilisation, sa gestion des sessions, son intégration avec LangChain et la rapidité de développement qu'il permet. Il facilite aussi le déploiement de l'application.

\subsubsection{Bibliothèques Python}

Nous avons mobilisé plusieurs bibliothèques Python au cours des phases 5, 6 et 7. Le tableau \ref{tab:bibliotheques} en présente la liste et le rôle.

\begin{table}[H]
\centering
\caption{Bibliothèques Python utilisées}
\label{tab:bibliotheques}
\begin{tabular}{|L{5cm}|L{3cm}|L{5cm}|}
\hline
\textbf{Bibliothèque} & \textbf{Version} & \textbf{Rôle} \\
\hline
\texttt{neo4j} & 5.28 & Pilote de connexion à Neo4j \\
\hline
\texttt{langchain}, \texttt{langchain-core} & 0.3 & Consignes et appel des modèles \\
\hline
\texttt{langchain-openai} & 0.3 & Accès à l'API d'Ollama \\
\hline
\texttt{langchain-google-genai} & 2.1 & Accès à Gemini \\
\hline
\texttt{streamlit} & 1.64 & Interface web \\
\hline
\texttt{plotly} & 7.1 & Visualisation du graphe \\
\hline
\texttt{rdflib} & 7.6 & Lecture du fichier Turtle \\
\hline
\texttt{pytest} & 9.1 & Tests logiciels \\
\hline
\texttt{pyngrok} & Sur Colab & Ouverture du tunnel ngrok \\
\hline
\end{tabular}
\end{table}

Ces bibliothèques couvrent les besoins de la recherche, notamment la connexion à la base, l'orchestration du système RAG, le développement de l'interface, le chargement du modèle, la manipulation de l'ontologie et le traitement des données.

% ------------------------------------------------------------
% II.1.4. Tableau récapitulatif du matériel
% ------------------------------------------------------------

\subsection{Tableau récapitulatif du matériel}
\label{subsec:recap_materiel}

Le tableau \ref{tab:recap_materiel} rassemble l'ensemble du matériel mobilisé.

\begin{table}[H]
\centering
\caption{Récapitulatif du matériel}
\label{tab:recap_materiel}
\begin{tabular}{|L{3cm}|L{4cm}|L{6cm}|}
\hline
\textbf{Catégorie} & \textbf{Élément} & \textbf{Spécification} \\
\hline
\multirow{3}{*}{Corpus} & Documents scientifiques & 35 rapports, articles, thèses \\
\cline{2-3}
& Base burkinaherbiermed.org & Plusieurs centaines de plantes \\
\cline{2-3}
& Base INERA & > 7 000 espèces géolocalisées \\
\hline
\multirow{5}{*}{Ontologies} & WATRIMed & OWL, Afrique de l'Ouest \\
\cline{2-3}
& ontoMEDTRAD & OWL + iconique, Afrique de l'Ouest \\
\cline{2-3}
& ATHMed & Cadre proposé, Ghana \\
\cline{2-3}
& Deep Ontologies & OWL, Afrique \\
\cline{2-3}
& FLOPO, PO & OWL, botanique \\
\hline
Modèles de langage & gpt-6-luna ; Llama 3 8B & OpenWork (extraction) ; Ollama sur Colab et ngrok (RAG) \\
\hline
Plateforme & Google Colab & GPU T4 \\
\hline
Base de connaissances & Neo4j Desktop & Orientée graphe, Cypher \\
\hline
Orchestration RAG & LangChain & Consignes et appel des modèles \\
\hline
Interface web & Streamlit & Chatbot, recherche, fiches \\
\hline
Bibliothèques & Python & neo4j, langchain, streamlit, rdflib, etc. \\
\hline
\end{tabular}
\end{table}

% ------------------------------------------------------------
% II.1.5. Tableau de correspondance Matériel - Phases
% ------------------------------------------------------------

\subsection{Tableau de correspondance Matériel - Phases}
\label{subsec:correspondance_materiel}

Le tableau \ref{tab:correspondance_materiel} met en regard le matériel et les phases méthodologiques, et indique à quel moment et dans quel but chaque ressource entre en jeu.

\begin{table}[H]
\centering
\caption{Correspondance entre le matériel et les phases}
\label{tab:correspondance_materiel}
\begin{tabular}{|L{4cm}|L{3cm}|L{6cm}|}
\hline
\textbf{Matériel} & \textbf{Phase} & \textbf{Utilisation} \\
\hline
Corpus documentaire & Phase 3 & Constitution du corpus \\
\hline
Ontologies candidates & Phase 1 & Comparaison et sélection \\
\hline
gpt-6-luna & Phase 5 & Extraction \\
\hline
Llama 3 8B & Phase 7, Phase 8 & Génération RAG, évaluation \\
\hline
Google Colab & Phase 7, Phase 8 & Exécution de Llama 3 8B \\
\hline
Neo4j Desktop & Phase 6 & Base de connaissances \\
\hline
LangChain & Phase 7 & Orchestration RAG \\
\hline
Streamlit & Phase 7 & Interface web \\
\hline
Bibliothèques Python & Phase 5, 6, 7 & Développement \\
\hline
\end{tabular}
\end{table}

% ============================================================
% II.2. MÉTHODOLOGIE
% ============================================================

\section{Méthodologie}
\label{sec:methodologie}

Nous avons organisé la démarche en huit phases opérationnelles. Nous avons comparé les ontologies, défini les critères de recherche et constitué le corpus. Nous avons adapté l'ontologie, extrait les connaissances, construit la base de connaissances, développé le système d'accès et l'application web, puis évalué l'ensemble. Nous avons choisi l'ontologie avant de constituer le corpus afin de fixer le cadre conceptuel et de guider l'extraction.

% ------------------------------------------------------------
% II.2.1. Phase 1 : Identification et étude comparative des ontologies
% ------------------------------------------------------------

\subsection{Phase 1 : Identification et étude comparative des ontologies}
\label{subsec:phase1}

\textbf{Objectif} : Sélectionner l'ontologie la plus pertinente pour notre contexte.

Nous avons comparé les ontologies en trois temps. Nous avons défini une grille de critères, attribué une note à chacune, puis retenu celle qui a obtenu le meilleur score. Ce choix a fixé le modèle conceptuel des étapes suivantes.

\subsubsection{Définition de la grille de critères}

Nous avons retenu huit critères adaptés à nos besoins : couverture du domaine, richesse sémantique, disponibilité, documentation, interopérabilité, adaptabilité au contexte burkinabè, traçabilité et prise en compte des savoirs locaux. Le tableau \ref{tab:grille_criteres} présente ces critères.

\begin{table}[H]
\centering
\caption{Grille de critères de comparaison}
\label{tab:grille_criteres}
\begin{tabular}{|C{0.5cm}|L{3cm}|L{6.6cm}|C{2.5cm}|}
\hline
\textbf{N°} & \textbf{Critère} & \textbf{Description} & \textbf{Pondération} \\
\hline
1 & Couverture du domaine & Capacité à représenter plantes, usages, propriétés, maladies & Élevée (3) \\
\hline
2 & Richesse sémantique & Nombre et qualité des classes, relations, axiomes & Élevée (3) \\
\hline
3 & Disponibilité & Ontologie accessible et téléchargeable & Élevée (3) \\
\hline
4 & Documentation & Qualité de la documentation et des exemples & Moyenne (2) \\
\hline
5 & Interopérabilité & Compatibilité avec OWL, RDF, SPARQL & Moyenne (2) \\
\hline
6 & Adaptabilité & Facilité d'adaptation au contexte burkinabè & Élevée (3) \\
\hline
7 & Traçabilité & Capacité à associer les connaissances aux sources & Élevée (3) \\
\hline
8 & Prise en compte des savoirs locaux & Intégration des noms vernaculaires, recettes multi-plantes & Élevée (3) \\
\hline
\end{tabular}
\end{table}

Nous avons attribué à chaque critère une pondération proportionnelle à son importance. Les critères fortement pondérés (3) étaient déterminants pour représenter le contexte burkinabè, tandis que les critères moyennement pondérés (2) autorisaient davantage de compromis. Nous avons calculé le score global en multipliant chaque note par la pondération correspondante, en additionnant les produits et en divisant cette somme par la somme des pondérations. Le résultat est rapporté sur 5.

\subsubsection{Notation des ontologies}

Nous avons noté chaque ontologie de 1 (très faible) à 5 (très élevé) pour chacun des critères. Le tableau \ref{tab:notation_ontologies} présente les résultats.

\begin{table}[H]
\centering
\caption{Notation des ontologies candidates}
\label{tab:notation_ontologies}
\begin{tabular}{|L{3.5cm}|C{0.8cm}|C{0.8cm}|C{0.8cm}|C{0.8cm}|C{0.8cm}|C{0.8cm}|C{0.8cm}|C{0.8cm}|C{1.5cm}|}
\hline
\textbf{Ontologie} & \textbf{C1} & \textbf{C2} & \textbf{C3} & \textbf{C4} & \textbf{C5} & \textbf{C6} & \textbf{C7} & \textbf{C8} & \textbf{Score} \\
\hline
WATRIMed & 4 & 4 & 2 & 3 & 5 & 4 & 3 & 5 & 3.7 \\
\hline
ontoMEDTRAD & 5 & 5 & 5 & 4 & 4 & 5 & 4 & 5 & 4.7 \\
\hline
ATHMed & 4 & 4 & 3 & 3 & 4 & 3 & 2 & 3 & 3.2 \\
\hline
Deep Ontologies & 3 & 4 & 2 & 2 & 3 & 3 & 2 & 4 & 2.9 \\
\hline
FLOPO / PO & 2 & 3 & 5 & 4 & 5 & 2 & 1 & 1 & 2.7 \\
\hline
\end{tabular}
\end{table}

Nous avons analysé les notes afin de comparer la couverture du domaine et la prise en compte des savoirs locaux par les ontologies africaines. Nous avons également examiné la disponibilité et l'interopérabilité des ontologies botaniques, ainsi que les différences entre WATRIMed et ontoMEDTRAD concernant la richesse du modèle et l'accessibilité de la documentation.

Nous avons accordé une attention particulière à la disponibilité (C3) et à la documentation (C4). Nous avons évalué les ressources effectivement accessibles. WATRIMed est décrite par deux articles de congrès, dont l'un compte deux pages, mais son fichier d'ontologie n'est pas téléchargeable. La thèse consacrée à ontoMEDTRAD fournit la liste des classes et des propriétés, les métriques ainsi que l'origine des concepts.

\subsubsection{Sélection de l'ontologie}

WATRIMed \cite{watrimed2019, watrimed_design_2019} pourrait se distinguer par la richesse de son modèle documenté, qui compte 472 concepts et 75 propriétés, ainsi que par son alignement sur l'ontologie fondationnelle BioTopLite2. Toutefois, les deux publications n'en décrivent que les classes supérieures et les relations principales, et le fichier de l'ontologie n'est pas disponible pour une réutilisation directe.

Nous avons retenu ontoMEDTRAD \cite{kouame_systeme_2018}, qui a obtenu le meilleur score. Cette ontologie couvre l'Afrique de l'Ouest, représente les recettes multi-plantes, leurs ingrédients, les parties utilisées et les modes de préparation, et documente l'origine de ses concepts à partir de quinze ressources terminologiques et ontologiques. Sa thèse est accessible. Nous avons aussi pris en compte ses limites : sa validation globale reste annoncée comme perspective et son jeu de données comprend 22 plantes et 31 remèdes antipaludéens \cite{kouame_visual_2020}.

L'ontologie retenue fournira le cadre conceptuel pour extraire et organiser les connaissances. Elle guidera la constitution du corpus et son adaptation au contexte burkinabè.

\textbf{Livrable} : Ontologie ontoMEDTRAD sélectionnée.

% ------------------------------------------------------------
% II.2.2. Phase 2 : Définition des critères de recherche et de sélection
% ------------------------------------------------------------

\subsection{Phase 2 : Définition des critères de recherche et de sélection des documents}
\label{subsec:phase2}

\textbf{Objectif} : Définir les critères qui guideront la recherche et la sélection des documents scientifiques.

Nous avons défini deux catégories de critères, l'une pour rechercher les informations et l'autre pour sélectionner les documents. Elles visaient à assurer la pertinence, la qualité et la cohérence des sources retenues.

\subsubsection{Critères de recherche}

Nous avons retenu les critères suivants pour guider la recherche :

\begin{itemize}
\item \textbf{Les plantes} : espèces végétales documentées au Burkina Faso, avec identification botanique (nom scientifique, famille, synonymes).
\item \textbf{Les parties utilisées} : feuilles, racines, écorces, fruits, graines, fleurs, tiges, etc.
\item \textbf{Les pathologies et usages médicinaux} : paludisme, infections respiratoires, digestives, diabète, hypertension, etc.
\item \textbf{Les modes de préparation} : décoction, macération, infusion, poudre, extrait alcoolique, etc.
\item \textbf{Les niveaux de preuve} : savoirs traditionnels, études in vitro, études in vivo, essais cliniques.
\item \textbf{Les propriétés pharmacologiques} : antimicrobiennes, antipaludiques, anti-inflammatoires, antioxydantes, toxicité, etc.
\item \textbf{La composition chimique} : alcaloïdes, flavonoïdes, tanins, saponines, huiles essentielles, etc.
\end{itemize}

Ils couvrent les principales dimensions des connaissances sur les plantes médicinales et ont permis de cibler les informations utiles à l'étude.

\subsubsection{Critères de sélection des documents}

Nous avons sélectionné les documents selon les critères suivants :

\begin{itemize}
\item \textbf{Pertinence} : le document traite des plantes médicinales du Burkina Faso.
\item \textbf{Identification de la plante} : le nom scientifique est précisé.
\item \textbf{Disponibilité d'une référence scientifique} : le document est publié dans une revue à comité de lecture, une thèse, un mémoire ou un rapport de recherche.
\item \textbf{Fiabilité} : les sources sont vérifiables et citées.
\item \textbf{Actualité} : les documents récents sont privilégiés, sans exclure les travaux fondateurs.
\end{itemize}

Ces critères nous ont permis d'écarter les documents non scientifiques ou non pertinents et de constituer un corpus cohérent pour l'étude.

\textbf{Livrable} : grille de critères de recherche et de sélection.

% ------------------------------------------------------------
% II.2.3. Phase 3 : Constitution du corpus scientifique
% ------------------------------------------------------------

\subsection{Phase 3 : Constitution du corpus scientifique}
\label{subsec:phase3}

\textbf{Objectif} : Rassembler un corpus documentaire fiable et pertinent sur les plantes médicinales du Burkina Faso.

Nous avons constitué le corpus en quatre étapes. Nous avons recherché les documents dans Google Scholar, PubMed et ResearchGate, puis consulté burkinaherbiermed.org et les données de l'INERA. Nous avons inclus les documents portant sur les plantes médicinales du Burkina Faso, précisant le nom scientifique des plantes et publiés ou validés par une institution scientifique avec des sources vérifiables. Nous avons écarté les documents non scientifiques, ceux sans identification botanique et les documents anciens sans valeur de référence. Nous avons classé les documents par type, plante, pathologie et région et les avons enregistrés dans un répertoire structuré. Le corpus comprend 35 documents scientifiques et les données de burkinaherbiermed.org et de l'INERA.

\textbf{Livrable} : un corpus de 35 documents scientifiques, complété par les données de burkinaherbiermed.org et de l'INERA.

% ------------------------------------------------------------
% II.2.4. Phase 4 : Adaptation de l'ontologie sélectionnée
% ------------------------------------------------------------

\subsection{Phase 4 : Adaptation de l'ontologie sélectionnée}
\label{subsec:phase4}

\textbf{Objectif} : Adapter le modèle conceptuel de l'ontologie sélectionnée au contexte burkinabè.

Nous avons examiné les concepts existants, ajouté des classes et des propriétés, modifié certaines relations et documenté les changements.

\subsubsection{Analyse des concepts existants}

Nous avons étudié la structure de l'ontologie ontoMEDTRAD telle qu'elle est décrite dans la thèse \cite{kouame_systeme_2018}. L'ontologie compte 144 concepts, 76 propriétés associatives, 16 propriétés de données, 789 individus et 2 734 axiomes. Elle est formalisée en logique de description ALCHOIQ, éditée sous Protégé~5 et vérifiée par les raisonneurs HermiT et Pellet.

Le fichier OWL de l'ontologie n'étant pas téléchargeable publiquement, nous avons reconstitué sa structure à partir de la thèse. Nous nous sommes appuyés sur les captures d'écran de Protégé (figures 5.15 à 5.20) et les tableaux récapitulatifs des concepts, propriétés et métriques (tableaux 5.5, 5.21 et 5.22). Cette reconstitution a servi de base à l'adaptation présentée au chapitre \ref{chap:conception_developpement}.

Nous avons reconstitué la structure conceptuelle de l'ontologie, ses classes, propriétés et relations, mais pas ses instances d'origine. Nous n'avons pas réutilisé les 789 individus d'ontoMEDTRAD, liés au contexte et aux sources des travaux d'origine. Cette phase a produit le modèle conceptuel adapté ; le peuplement par les individus extraits de la littérature scientifique est intervenu en phase 5.

Ses concepts principaux décrivent le parcours de soin, du praticien et du patient aux symptômes, aux maladies, aux ressources médicinales, à l'examen clinique, aux recettes, aux ingrédients, aux parties utilisées, aux préparations et aux remèdes. L'ontologie distingue cinq types de plantes et neuf parties utilisées. Elle comprend aussi une classe pour les références scientifiques associées aux ressources, aux recettes et aux remèdes. Les concepts secondaires couvrent les contre-indications, les propriétés thérapeutiques, la toxicité, l'étiologie et la classification des maladies selon la CIM-10. Les relations de base expriment la subsomption et la relation partie-tout. Soixante-seize propriétés associatives complètent le modèle, dont la hiérarchie \textit{aPourForme}, utilisée pour l'iconisation des plantes.

Nous avons ensuite évalué ce modèle au regard du contexte burkinabè. Il couvre les concepts de base et les raisonneurs testés n'y ont signalé aucune incohérence, mais il nécessitait des adaptations pour représenter les noms vernaculaires en langues locales, les maladies spécifiques, les niveaux de preuve, les dosages et la traçabilité détaillée des sources. Sa validation globale reste une perspective, et son jeu de données se limite à 22 plantes et 31 remèdes antipaludéens.

\subsubsection{Ajout de classes}

La comparaison des fichiers du modèle conceptuel adapté et du schéma conservé dénombre 40 classes ajoutées. Le modèle adapté déclare également 25 propriétés objet au total, dont 20 absentes du schéma conservé, et 20 propriétés de données, dont 15 ajoutées. Les classes ajoutées couvrent notamment les noms vernaculaires, les usages thérapeutiques, les niveaux de preuve, les sources, les préparations, des parties végétales complémentaires, les composés chimiques, les activités pharmacologiques et les effets toxicologiques. Ces ajouts définissent le modèle conceptuel ; les individus et données extraits sont intégrés lors de la phase suivante.

Ces ajouts structurent des connaissances locales et scientifiques à différents niveaux. Ils ne préchargent pas les valeurs, comme les noms, pathologies, régions ou composés ; celles-ci proviennent des sources du corpus et sont enregistrées dans la base avec leur provenance.

\subsubsection{Ajout de propriétés}

Nous avons ajouté des propriétés de données pour les noms locaux, le niveau de preuve, le type d'étude, les références, le dosage et la toxicité, en complément de propriétés textuelles pour les noms vernaculaires et les descriptions d'activités ou d'effets. Des propriétés objet relient également l'usage à sa maladie cible, à la partie végétale utilisée, à la préparation et à la source ; d'autres décrivent les associations avec des composés, activités pharmacologiques, risques et localités. Les noms techniques exacts sont ceux déclarés dans la version courante du fichier Turtle.

L'enrichissement ainsi obtenu saisit des informations fines (noms vernaculaires, coordonnées, niveaux de preuve, références précises) que l'ontologie d'origine ne documentait pas pour le Burkina Faso.

\subsubsection{Modification de relations}

Nous avons structuré les liens entre plante, usage, maladie, partie utilisée et préparation, puis généralisé la traçabilité documentaire. La propriété objet \textit{estDocumenteePar} relie les ressources du graphe à une source scientifique; des sous-propriétés spécialisées permettent de préciser si la source documente une plante, un usage, un nom vernaculaire, une activité, un composé ou un effet toxicologique. \textit{aPourNiveauDePreuve} rattache un usage au niveau déclaré par la source. Ces liens ne signifient pas que les usages ou activités constituent des recommandations thérapeutiques ni une validation clinique.

Ces modifications assurent la traçabilité des connaissances et permettent d'exprimer des relations que le modèle d'origine, fondé sur la relation « est un » et sur la relation partie-tout, ne couvrait pas.

\subsubsection{Documentation des adaptations}

Nous avons consigné toutes les modifications dans un journal : chaque ajout est justifié, chaque choix est explicité, et les classes, propriétés et relations ajoutées, modifiées ou supprimées y sont décrites une à une.

Cette documentation n'est pas une formalité. Elle rend la démarche reproductible par un tiers et facilite les évolutions ultérieures de l'ontologie.

\textbf{Livrable} : Le modèle conceptuel adapté de l'ontologie, sans les données instanciées.

% ------------------------------------------------------------
% II.2.5. Phase 5 : Enrichissement de l'ontologie
% ------------------------------------------------------------

\subsection{Phase 5 : Enrichissement des connaissances à partir de la littérature scientifique}
\label{subsec:phase5}

\textbf{Objectif} : Extraire les connaissances du corpus pour alimenter la base organisée selon l'ontologie.

À partir du modèle conceptuel adapté en phase 4, nous avons enrichi l'ontologie avec des individus et des données extraits des sources scientifiques. Nous avons préparé ces données structurées et sourcées pour leur exploitation dans la base orientée graphe Neo4j.

Nous avons conçu les consignes d'extraction, procédé à l'extraction assistée par modèle de langue, contrôlé les données obtenues, puis structuré et préparé ces données pour leur intégration dans la base.

\subsubsection{Conception des prompts d'extraction}

Nos consignes d'extraction ciblaient les noms scientifiques et vernaculaires, la famille botanique, les parties utilisées, les usages, les préparations, les activités pharmacologiques, le niveau de preuve et la source. Le chapitre \ref{chap:conception_developpement} présente un exemple de consigne, avec le format JSON et les champs demandés.

Nous avons employé une même consigne pour les documents traités. Les champs absents reçoivent la valeur \texttt{null} ou une liste vide, et aucune valeur n'est complétée par déduction.

\subsubsection{Exécution de l'extraction}

Nous avons soumis les documents à gpt-6-luna dans l'environnement OpenWork, un par un, et enregistré les résultats en JSON. L'extraction a porté sur dix des 35 documents du corpus, principalement des enquêtes ethnobotaniques ; les 25 autres n'ont pas été extraits.

\subsubsection{Contrôle de la qualité des extractions}

Nous avons repris le niveau de preuve indiqué par la source lorsqu'il était explicite, qu'il s'agisse d'un savoir traditionnel, d'une étude \textit{in vitro}, d'une étude \textit{in vivo} ou d'un essai clinique. Nous n'avons pas évalué indépendamment la solidité des faits et avons laissé le champ vide lorsque la source ne précisait pas ce niveau. Un expert externe n'a pas relu les extractions. Nous avons contrôlé leur structure, puis comparé une partie des faits aux documents d'origine à la phase 8.

\subsubsection{Structuration des données}

Nous avons organisé les données selon le modèle de l'ontologie dans un classeur d'individus, par type de connaissance (sources, plantes, noms vernaculaires, usages, préparations, maladies, activités, composés, toxicologie et localités), puis les avons transformées en individus au format Turtle.

\subsubsection{Préparation de l'intégration dans la base de connaissances}

Le fichier OntoPlantMedBF-peuplee.ttl contient des individus extraits de dix documents scientifiques, en plus du schéma adapté. Ce fichier constitue l'ontologie enrichie importée dans Neo4j à la phase 6. Les données de burkinaherbiermed.org et de l'INERA, consultées lors de la constitution du corpus, n'y ont pas été intégrées.

\textbf{Livrable} : OntoPlantMedBF, ontologie enrichie avec les données extraites et leurs sources.

% ------------------------------------------------------------
% II.2.6. Phase 6 : Construction et interrogation de la base
% ------------------------------------------------------------

\subsection{Phase 6 : Construction et interrogation de la base de connaissances}
\label{subsec:phase6}

\textbf{Objectif} : Importer et exploiter dans Neo4j l'ontologie réajustée \textit{OntoPlantMedBF}.

Nous avons importé dans Neo4j l'ontologie réajustée \textit{OntoPlantMedBF} et ses données, en les représentant sous forme de nœuds, propriétés et relations du graphe. Nous avons vérifié la correspondance avec le modèle ontologique et testé les requêtes de la base.

\subsubsection{Import et ajustement du modèle dans le graphe}

Nous avons importé les classes et propriétés d'OntoPlantMedBF dans Neo4j sous forme de labels, propriétés et relations. Les principaux nœuds représentent les plantes, maladies, usages, parties, préparations, sources, composés chimiques et régions. Les relations du graphe reprennent les liens du modèle, avec des noms adaptés aux conventions Neo4j, et les propriétés conservent notamment les noms, niveaux de preuve, DOI et coordonnées géographiques.

Nous avons ainsi importé dans le graphe la structure réajustée d'OntoPlantMedBF et ses données. Les entités sont représentées par des nœuds, leurs liens par des relations et leurs caractéristiques par des propriétés, conformément au modèle ontologique.

\subsubsection{Import des données}

Nous avons chargé les données préparées, créé ou rapproché les nœuds et relations au moyen de requêtes Cypher, puis vérifié l'intégrité et la cohérence de l'import avec le modèle ontologique.

Après l'import, nous avons contrôlé que les données, les relations et les propriétés avaient été correctement reprises.

\subsubsection{Requêtes de test}

Nous avons testé quatre familles de requêtes : les requêtes simples sur les plantes associées à une maladie, les requêtes complexes sur les recettes multi-plantes, les requêtes de traçabilité et les requêtes spatiales par région.

Ces vérifications nous ont permis de contrôler le fonctionnement de la base et les résultats renvoyés par les requêtes.

\textbf{Livrable} : Une base Neo4j peuplée et interrogeable, complémentaire d'OntoPlantMedBF.

% ------------------------------------------------------------
% II.2.7. Phase 7 : Exploitation par l'IA et développement web
% ------------------------------------------------------------

\subsection{Phase 7 : Exploitation par IA et développement de l'application web}
\label{subsec:phase7}

\textbf{Objectif} : Développer le système d'accès intelligent et l'interface utilisateur.

Nous avons mis en place le système RAG en définissant son architecture, en construisant la chaîne de traitement et en intégrant ses composants.

\subsubsection{Architecture du système}

L'architecture relie l'utilisateur à Streamlit, LangChain, Neo4j et Llama 3 jusqu'à la réponse. Elle associe la recherche dans la base à la génération par le modèle de langue et conserve le lien entre les informations et leurs sources. Le chapitre \ref{chap:conception_developpement} décrit la chaîne de traitement, de l'interprétation de la question à l'affichage des références.

\subsubsection{Intégration des composants}

Nous avons connecté l'application à Neo4j et au modèle de langue, puis configuré les consignes de génération des réponses. Cette intégration a réuni les composants du système RAG.

\subsubsection{Développement de l'application web}

Nous avons conçu l'application autour d'une page d'accueil avec chatbot, d'une recherche par plante, d'une recherche par pathologie et de fiches descriptives. Nous y avons intégré le système RAG et la gestion de l'historique. L'interface permet d'interroger la base en langage naturel et de recevoir des réponses reliées à leurs sources. Le chapitre \ref{chap:conception_developpement} en décrit la conception et les composants.

\subsubsection{Tests et validation}

Nous avons mené des tests fonctionnels pour vérifier le fonctionnement de l'application et des tests utilisateurs pour examiner sa facilité d'utilisation et la pertinence des réponses.

\textbf{Livrable} : Un système RAG fonctionnel et une application web fonctionnelle.

% ------------------------------------------------------------
% II.2.8. Phase 8 : Évaluation et validation
% ------------------------------------------------------------

\subsection{Phase 8 : Évaluation et validation}
\label{subsec:phase8}

\textbf{Objectif} : Évaluer la qualité de l'ontologie, du système RAG et de l'application web.

Nous avons évalué trois composantes : l'ontologie, le système RAG et l'application web.

\subsubsection{Évaluation de l'ontologie}

Nous avons évalué l'ontologie selon les quatre critères présentés dans le tableau \ref{tab:eval_ontologie}.

\begin{table}[H]
\centering
\caption{Critères d'évaluation de l'ontologie}
\label{tab:eval_ontologie}
\begin{tabular}{|L{3cm}|L{6cm}|L{4cm}|}
\hline
\textbf{Critère} & \textbf{Description} & \textbf{Méthode} \\
\hline
Couverture & Capacité à représenter les connaissances du corpus & Vérification manuelle \\
\hline
Cohérence & Absence de contradictions & Requêtes de test \\
\hline
Traçabilité & Association des connaissances aux sources & Vérification des relations \\
\hline
Extensibilité & Facilité d'ajout de nouvelles connaissances & Test d'ajout \\
\hline
\end{tabular}
\end{table}

\subsubsection{Évaluation du système RAG}

Nous avons évalué le système RAG selon les quatre critères présentés dans le tableau \ref{tab:eval_rag}.

\begin{table}[H]
\centering
\caption{Critères d'évaluation du système RAG}
\label{tab:eval_rag}
\begin{tabular}{|L{3cm}|L{6cm}|L{4cm}|}
\hline
\textbf{Critère} & \textbf{Description} & \textbf{Métrique} \\
\hline
Pertinence & Les réponses répondent à la question & Précision \\
\hline
Fidélité & Les réponses sont fidèles aux sources & Rappel \\
\hline
Exactitude & Les réponses sont correctes & F1-score \\
\hline
Traçabilité & Les sources sont affichées & Taux de couverture \\
\hline
\end{tabular}
\end{table}

Nous avons constitué un jeu de 50 questions de référence portant sur les plantes, les usages et leurs sources, puis comparé les réponses du système aux réponses attendues.

\subsubsection{Évaluation de l'application web}

Nous avons évalué l'application web selon les quatre critères présentés dans le tableau \ref{tab:eval_web}.

\begin{table}[H]
\centering
\caption{Critères d'évaluation de l'application web}
\label{tab:eval_web}
\begin{tabular}{|L{3cm}|L{6cm}|L{4cm}|}
\hline
\textbf{Critère} & \textbf{Description} & \textbf{Méthode} \\
\hline
Facilité d'utilisation & L'interface est intuitive & Questionnaire \\
\hline
Compréhension & Les réponses sont claires & Questionnaire \\
\hline
Utilité & L'application répond aux besoins & Questionnaire \\
\hline
Performance & Temps de réponse & Mesure \\
\hline
\end{tabular}
\end{table}

Nous avons sollicité des étudiants, des chercheurs et des tradipraticiens pour cette évaluation.

\textbf{Livrable} : Rapport d'évaluation et de validation.

% ============================================================
% II.3. SYNTHÈSE ET TRANSITION
% ============================================================

\section{Synthèse et transition}
\label{sec:synthese}

\subsection{Tableau de correspondance Phases - Livrables}

Le tableau \ref{tab:correspondance_phases} récapitule les objectifs et les livrables des huit phases.

\begin{table}[H]
\centering
\caption{Correspondance entre les phases, les objectifs et les livrables}
\label{tab:correspondance_phases}
\begin{tabular}{|L{2.2cm}|L{4.5cm}|L{6.8cm}|}
\hline
\textbf{Phase} & \textbf{Objectif} & \textbf{Livrable} \\
\hline
Phase 1 & Sélectionner l'ontologie la plus pertinente & Ontologie ontoMEDTRAD sélectionnée \\
\hline
Phase 2 & Définir les critères de recherche et de sélection & Grille de critères \\
\hline
Phase 3 & Rassembler un corpus fiable et pertinent & Corpus de 35 documents \\
\hline
Phase 4 & Adapter le modèle conceptuel de l'ontologie & Modèle conceptuel adapté \\
\hline
Phase 5 & Extraire les connaissances et enrichir l'ontologie & Ontologie enrichie en données \\
\hline
Phase 6 & Importer OntoPlantMedBF dans la base graphe & OntoPlantMedBF intégrée à Neo4j, base interrogeable \\
\hline
Phase 7 & Développer le système RAG et l'application web & Système RAG et application web \\
\hline
Phase 8 & Évaluer la qualité du système & Rapport d'évaluation \\
\hline
\end{tabular}
\end{table}

% ============================================================
% CONCLUSION DU CHAPITRE
% ============================================================

\section*{Conclusion du chapitre}
\addcontentsline{toc}{section}{Conclusion du chapitre}

Ce chapitre a exposé le matériel et la méthodologie de notre recherche. Le matériel réunit un corpus de 35 documents scientifiques, les données de burkinaherbiermed.org et de l'INERA, cinq ontologies candidates et un environnement technique bâti sur gpt-6-luna pour l'extraction, Llama 3 8B sur Google Colab pour le RAG, Neo4j Desktop, LangChain et Streamlit.

Nous avons suivi huit phases : comparaison des ontologies, définition des critères de recherche et de sélection, constitution du corpus, adaptation de l'ontologie, extraction des connaissances, construction et interrogation de la base, développement du système d'accès et de l'application, puis évaluation. Nous avons décrit les actions menées et les livrables associés.

Nous avons choisi l'ontologie avant de constituer le corpus afin d'assurer la cohérence entre le modèle conceptuel et les données collectées, et de cibler les informations pertinentes dès la collecte.

Nous signalons deux limites. Les niveaux de preuve associés aux connaissances proviennent des sources et n'ont pas fait l'objet d'une évaluation indépendante. De plus, un expert externe n'a pas validé les extractions. Ces limites sont à prendre en compte dans les travaux ultérieurs. Nous avons néanmoins conservé le lien entre chaque information et sa source.

Le chapitre \ref{chap:resultats_discussion} présente les résultats de l'évaluation menée en phase 8.

Le chapitre suivant présente la conception et le développement du système, notamment l'ontologie adaptée, la base de connaissances, le système RAG et l'application web.
).

\phantomsection
\section*{Introduction du chapitre}

Après avoir posé le cadre conceptuel et exposé l'état de l'art, nous présentons le matériel mobilisé et la démarche méthodologique suivie. Ce chapitre fait le lien entre la réflexion théorique et la mise en œuvre technique.

Nous présentons le matériel mobilisé, notamment le corpus documentaire, les ontologies candidates et l'environnement technique. Nous décrivons les huit phases de la démarche suivie, de l'étude comparative des ontologies à l'évaluation du système, puis en proposons une synthèse.

% ============================================================
% II.1. MATÉRIEL
% ============================================================

\section{Matériel}
\label{sec:materiel}

Le matériel comprend le corpus documentaire, les ontologies candidates et l'environnement technique. Les deux tableaux qui suivent récapitulent ces ressources, réparties entre documentation, modèles conceptuels et outils techniques.

% ------------------------------------------------------------
% II.1.1. Corpus documentaire
% ------------------------------------------------------------

\subsection{Corpus documentaire}
\label{subsec:corpus}

Nous avons choisi l'ontologie avant de constituer le corpus afin de déterminer les notions et les relations qui guideraient l'extraction. Nous avons réuni les publications scientifiques et les données institutionnelles consultées sur burkinaherbiermed.org et auprès de l'INERA.

La figure \ref{fig:flux_corpus} résume le processus de constitution du corpus, du choix de l'ontologie à la sélection des documents.

\begin{figure}[H]
\centering
\includegraphics[width=0.8\textwidth]{flux_corpus.png}
\caption{Constitution du corpus documentaire}
\label{fig:flux_corpus}
\end{figure}

\subsubsection{Sources scientifiques publiées}

Notre corpus scientifique comprend 35 rapports de recherche, articles, thèses et mémoires provenant de revues, d'universités et d'institutions de recherche burkinabè. Ces travaux ethnobotaniques, ethnopharmacologiques et phytochimiques documentent les noms scientifiques et vernaculaires, les parties utilisées, les pathologies traitées, les modes de préparation et, pour les études qui les rapportent, les dosages et les données de toxicité.

Nous avons retenu ces travaux pour leur qualité et leur pertinence au contexte burkinabè, des publications fondatrices aux plus récentes. Cette diversité d'espèces, de pathologies et de pratiques thérapeutiques soutient l'extraction. Quatre travaux illustrent les types de données mobilisés dans le corpus :
\begin{itemize}
\item Zerbo et al.~\cite{zerbo_plantes_2011} ont étudié les plantes médicinales chez les Sanan du Burkina Faso. Leur travail documente 85 espèces utilisées pour traiter 63 affections.
\item Ouédraogo et al.~\cite{ouedraogo_revue_2024} ont dressé une revue des plantes médicinales à potentiel aromatique du pays, recensant 42 espèces et en caractérisant les propriétés et les usages.
\item Tibiri et al.~\cite{tibiri2020} ont mené une enquête à l'échelle nationale sur les plantes utilisées contre les maladies hépatiques, documentant plus de 2 000 entrées de plantes correspondant à 212 espèces.
\item Ouédraogo et al.~\cite{ouedraogo2020markets} ont enquêté auprès de 30 commerçants de plantes médicinales dans six grandes villes, identifiant 104 espèces utilisées pour diverses pathologies.
\end{itemize}

Ces travaux fournissent des méthodes réutilisables et des observations de terrain recueillies au Burkina Faso.

\subsubsection{Données de l'Herbier National des Plantes Médicinales du Burkina Faso}

Le site burkinaherbiermed.org diffuse des données sur les plantes médicinales du pays, notamment des informations botaniques, des usages thérapeutiques et des références bibliographiques. Ces données proviennent du projet « Herbier avec chimiothèque », piloté par le Ministère de la Santé. Ce projet a pour objectif de constituer un herbier national dédié aux plantes médicinales et une chimiothèque de marqueurs chimiques. Il documente les plantes utilisées contre les dix pathologies prioritaires définies par le ministère, géolocalise les stations et conserve des spécimens. Ces informations alimentent une base nationale accessible par application mobile et site web. Le projet associe un herbier, une chimiothèque et une base de données ethnobotaniques et géographiques. Il complète les autres herbiers burkinabè, notamment le HNBU du CNRST/INERA, le Centre National de Semences Forestières et l'herbier de l'Université Joseph Ki-Zerbo.

\subsubsection{Données de l'INERA}

La base de données de l'Institut de l'Environnement et de Recherches Agricoles (INERA) comprend plus de 7 000 espèces végétales et les coordonnées géographiques des stations où elles ont été observées. Ces données complètent les informations ethnobotaniques par une dimension spatiale et permettent d'étudier la répartition des plantes au Burkina Faso. Le rapprochement des usages rapportés et de la distribution géographique peut faire apparaître des régularités régionales, qui restent à vérifier.

Le tableau \ref{tab:corpus} récapitule les sources du corpus.

\begin{table}[H]
\centering
\caption{Sources du corpus documentaire}
\label{tab:corpus}
\begin{tabular}{|L{4cm}|L{2.8cm}|L{3cm}|L{5cm}|}
\hline
\textbf{Source} & \textbf{Type} & \textbf{Volume estimé} & \textbf{Contenu principal} \\
\hline
Articles et thèses & Publications scientifiques & 35 documents & Études ethnobotaniques, pharmacologiques \\
\hline
burkinaherbiermed.org & Base de données en ligne & Plusieurs centaines de plantes & Données botaniques, usages, références \\
\hline
INERA & Base de données institutionnelle & > 7 000 espèces & Coordonnées géographiques, identification \\
\hline
\end{tabular}
\end{table}

Les documents du corpus relèvent de types d'étude différents, et chacun est exploitable pour les informations qu'il rapporte. Trois profils se distinguent. Les enquêtes ethnobotaniques et ethnopharmacologiques, menées auprès de tradipraticiens ou dans les marchés, rapportent les noms vernaculaires, les parties utilisées, les modes de préparation et les pathologies traitées, souvent plante par plante. Les études expérimentales, notamment les essais \textit{in vitro} et les travaux phytochimiques, mesurent des activités pharmacologiques, des concentrations, des niveaux de preuve et parfois les lieux de récolte. Elles renseignent moins sur les noms vernaculaires et les pratiques de préparation. Enfin, les revues de synthèse dressent des inventaires à grande échelle et servent surtout à repérer les sources primaires.

Nous avons appliqué trois règles à l'extraction. Chaque source ne renseignait que les champs qu'elle mesurait. Nous avons laissé vides les champs non couverts, sans les compléter par déduction. Nous n'avons pas retenu comme fait établi un usage traditionnel rapporté dans un article expérimental sans référence ethnobotanique. Enfin, nous n'avons pas additionné les nombres d'espèces d'une source à l'autre, car une revue peut reprendre les espèces documentées dans les enquêtes qu'elle cite. Le tableau \ref{tab:corpus} présente les sources document par document, sans agréger les volumes.

Ce corpus rassemble des études menées dans plusieurs régions du Burkina Faso, portant sur différentes pathologies et mobilisant des approches méthodologiques variées. Il a servi de base à l'extraction des connaissances et à l'enrichissement de l'ontologie. Le choix préalable de l'ontologie a guidé l'extraction et l'organisation des documents selon les concepts et les relations définis.

% ------------------------------------------------------------
% II.1.2. Ontologies candidates
% ------------------------------------------------------------

\subsection{Ontologies candidates}
\label{subsec:ontologies_candidates}

Nous avons retenu cinq ontologies candidates pour l'étude comparative. Le chapitre \ref{chap:cadre_theorique}, aux sections \ref{subsubsec:ontologies_internationales} et \ref{subsubsec:ontologies_africaines}, les présente. Elles couvrent différents domaines, notamment la botanique, la médecine traditionnelle, la pharmacologie et l'ethnomédecine. Le tableau \ref{tab:ontologies_candidates} en présente les caractéristiques.

Nous les avons comparées à l'aide de la grille de critères présentée en phase 1 (section \ref{subsec:phase1}). Nous avons sélectionné l'ontologie la plus adaptée à notre contexte avant de constituer le corpus, afin de fixer le modèle conceptuel qui a guidé l'extraction.

\begin{table}[H]
\centering
\caption{Ontologies candidates pour la comparaison}
\label{tab:ontologies_candidates}
\begin{tabular}{|C{0.5cm}|L{3cm}|L{3cm}|L{3cm}|L{2cm}|L{2.5cm}|}
\hline
\textbf{N°} & \textbf{Ontologie} & \textbf{Domaine} & \textbf{Zone géographique} & \textbf{Format} & \textbf{Source d'accès} \\
\hline
1 & WATRIMed & Médecine traditionnelle & Afrique de l'Ouest & OWL, graphe de connaissances & IOS Press (2 articles) \\
\hline
2 & ontoMEDTRAD & Médecine traditionnelle & Afrique de l'Ouest & OWL + iconique & Thèse A. Kouamé (UGB) \\
\hline
3 & ATHMed & Médecine traditionnelle & Afrique (Ghana) & Cadre conceptuel & Frontiers in AI \\
\hline
4 & Deep Ontologies & Médecine traditionnelle & Afrique & OWL (Protégé) & Articles ARIMA \\
\hline
5 & FLOPO, PO & Botanique & International & OWL & OBO Foundry \\
\hline
\end{tabular}
\end{table}

Chacune présente des acquis et des lacunes. Les ontologies botaniques (FLOPO, PO) décrivent bien la plante mais ne traitent pas des usages médicinaux. Les ontologies pharmaceutiques (DrOn) se resserrent sur les médicaments modernes. L'OTM, prometteuse, reste sans ancrage africain. Les initiatives africaines (WATRIMed, ontoMEDTRAD, ATHMed) s'intéressent en revanche à la médecine traditionnelle du terrain, avec ses dimensions culturelles et ses recettes multi-plantes, même si ATHMed en reste, elle, au stade d'un cadre proposé. La phase 1 en livre la comparaison détaillée.

% ------------------------------------------------------------
% II.1.3. Outils et environnement technique
% ------------------------------------------------------------

\subsection{Outils et environnement technique}
\label{subsec:environnement}

Les outils techniques occupent une place importante dans notre démarche. Nous décrivons les phases concernées et les raisons de leur choix.

La figure \ref{fig:dispositif_technique} résume l'environnement d'expérimentation, en précisant où chaque composant est exécuté et comment ils communiquent.

\begin{figure}[H]
\centering
\includegraphics[width=0.8\textwidth]{dispositif_technique.png}
\caption{Dispositif technique d'expérimentation}
\label{fig:dispositif_technique}
\end{figure}

\subsubsection{Modèles de langage}

Nous avons utilisé deux modèles, chacun pour une étape distincte. Pour l'extraction des connaissances (phase 5), nous avons utilisé \textbf{gpt-6-luna} d'OpenAI, accessible par l'environnement OpenWork. Pour la génération des réponses du système RAG (phase 7), nous avons utilisé \textbf{Llama 3 8B} \cite{meta2024llama3}, un modèle ouvert. Gemini (\textit{gemini-3.5-flash-lite}) nous a servi de point de comparaison lors de l'évaluation.

Le tableau \ref{tab:llama_acces} présente les modes d'accès à ces modèles.

\begin{table}[H]
\centering
\caption{Modèles de langage et modes d'accès}
\label{tab:llama_acces}
\begin{tabular}{|L{4cm}|L{4cm}|L{6cm}|}
\hline
\textbf{Usage} & \textbf{Modèle} & \textbf{Mode d'accès} \\
\hline
Extraction (phase 5) & gpt-6-luna & OpenWork \\
\hline
Génération RAG (phase 7) & Llama 3 8B & Ollama sur Google Colab, API exposée par ngrok \\
\hline
Comparaison (phase 8) & gemini-3.5-flash-lite & API de Google \\
\hline
\end{tabular}
\end{table}

Nous avons retenu Llama 3 pour la génération parce qu'il s'agit d'un modèle ouvert que nous pouvons exécuter sur une instance que nous contrôlons. La version 8B tient dans la mémoire d'un GPU accessible gratuitement sur Colab. Ollama sert le modèle et ngrok expose son API à l'application web, qui reste sur le poste local.

\subsubsection{Plateforme d'exécution : Google Colab}

Nous avons exécuté Llama 3 8B sur \textbf{Google Colab}, une plateforme de notebooks donnant accès à des GPU adaptés à l'inférence.

\textbf{Phases d'utilisation :} phases 7 et 8 pour la génération des réponses et l'évaluation du système RAG.

Cette plateforme nous a donné accès à un GPU sans nécessiter l'acquisition d'une infrastructure locale. La session reste cependant temporaire et l'adresse exposée par ngrok change à chaque lancement, ce qui convient à un prototype mais pas à un service pérenne.

\subsubsection{Base de connaissances : Neo4j Desktop}

Nous avons utilisé \textbf{Neo4j Desktop}, la version locale de Neo4j, pour créer, gérer et interroger la base de connaissances orientée graphe.

\textbf{Phase d'utilisation :} phase 6 pour la construction et l'interrogation de la base.

Le tableau \ref{tab:neo4j} présente les caractéristiques de Neo4j Desktop.

\begin{table}[H]
\centering
\caption{Caractéristiques de Neo4j Desktop}
\label{tab:neo4j}
\begin{tabular}{|L{5cm}|L{8cm}|}
\hline
\textbf{Caractéristique} & \textbf{Valeur} \\
\hline
Type & Base de données orientée graphe \\
\hline
Langage de requête & Cypher \\
\hline
Version & Desktop (locale) \\
\hline
Format d'import & CSV, JSON \\
\hline
Interface & Neo4j Browser, Bloom \\
\hline
\end{tabular}
\end{table}

Nous avons choisi Neo4j pour représenter les relations entre les plantes, les usages, les propriétés et les sources, interroger le graphe avec Cypher et assurer la traçabilité des connaissances. La base permet également d'ajouter des connaissances et de visualiser les liens entre les entités.

\subsubsection{Orchestration RAG : LangChain}

Nous avons utilisé \textbf{LangChain} pour définir les consignes et appeler les modèles de langage de façon uniforme. La recherche dans Neo4j passe par le pilote officiel et par nos propres requêtes.

\textbf{Phase d'utilisation :} phase 7 pour l'accès aux connaissances et le développement web.

Le tableau \ref{tab:langchain} présente les composants LangChain utilisés.

\begin{table}[H]
\centering
\caption{Composants LangChain utilisés}
\label{tab:langchain}
\begin{tabular}{|L{5cm}|L{8cm}|}
\hline
\textbf{Composant} & \textbf{Rôle} \\
\hline
PromptTemplate & Préparation des consignes de génération \\
\hline
ChatOpenAI & Appel de Llama 3 par l'API d'Ollama, compatible avec celle d'OpenAI \\
\hline
ChatGoogleGenerativeAI & Appel de Gemini pour la comparaison \\
\hline
\end{tabular}
\end{table}

Nous avons retenu LangChain parce qu'il permet de changer de modèle sans modifier le reste du système. Nous n'avons pas confié la sélection des données à une chaîne générant du Cypher : le chapitre \ref{chap:conception_developpement} justifie ce choix \cite{lewis2020rag}.

\subsubsection{Interface web : Streamlit}

Nous avons développé l'application web avec \textbf{Streamlit}, un framework Python qui permet de créer des applications web interactives.

\textbf{Phase d'utilisation :} phase 7 pour le développement de l'application web.

L'application propose un chatbot en langage naturel, des recherches par plante, pathologie ou région, des fiches descriptives et l'affichage des sources et des niveaux de preuve. Nous avons également prévu une fonction de visualisation du graphe de connaissances.

Nous avons retenu Streamlit pour sa simplicité d'utilisation, sa gestion des sessions, son intégration avec LangChain et la rapidité de développement qu'il permet. Il facilite aussi le déploiement de l'application.

\subsubsection{Bibliothèques Python}

Nous avons mobilisé plusieurs bibliothèques Python au cours des phases 5, 6 et 7. Le tableau \ref{tab:bibliotheques} en présente la liste et le rôle.

\begin{table}[H]
\centering
\caption{Bibliothèques Python utilisées}
\label{tab:bibliotheques}
\begin{tabular}{|L{5cm}|L{3cm}|L{5cm}|}
\hline
\textbf{Bibliothèque} & \textbf{Version} & \textbf{Rôle} \\
\hline
\texttt{neo4j} & 5.28 & Pilote de connexion à Neo4j \\
\hline
\texttt{langchain}, \texttt{langchain-core} & 0.3 & Consignes et appel des modèles \\
\hline
\texttt{langchain-openai} & 0.3 & Accès à l'API d'Ollama \\
\hline
\texttt{langchain-google-genai} & 2.1 & Accès à Gemini \\
\hline
\texttt{streamlit} & 1.64 & Interface web \\
\hline
\texttt{plotly} & 7.1 & Visualisation du graphe \\
\hline
\texttt{rdflib} & 7.6 & Lecture du fichier Turtle \\
\hline
\texttt{pytest} & 9.1 & Tests logiciels \\
\hline
\texttt{pyngrok} & Sur Colab & Ouverture du tunnel ngrok \\
\hline
\end{tabular}
\end{table}

Ces bibliothèques couvrent les besoins de la recherche, notamment la connexion à la base, l'orchestration du système RAG, le développement de l'interface, le chargement du modèle, la manipulation de l'ontologie et le traitement des données.

% ------------------------------------------------------------
% II.1.4. Tableau récapitulatif du matériel
% ------------------------------------------------------------

\subsection{Tableau récapitulatif du matériel}
\label{subsec:recap_materiel}

Le tableau \ref{tab:recap_materiel} rassemble l'ensemble du matériel mobilisé.

\begin{table}[H]
\centering
\caption{Récapitulatif du matériel}
\label{tab:recap_materiel}
\begin{tabular}{|L{3cm}|L{4cm}|L{6cm}|}
\hline
\textbf{Catégorie} & \textbf{Élément} & \textbf{Spécification} \\
\hline
\multirow{3}{*}{Corpus} & Documents scientifiques & 35 rapports, articles, thèses \\
\cline{2-3}
& Base burkinaherbiermed.org & Plusieurs centaines de plantes \\
\cline{2-3}
& Base INERA & > 7 000 espèces géolocalisées \\
\hline
\multirow{5}{*}{Ontologies} & WATRIMed & OWL, Afrique de l'Ouest \\
\cline{2-3}
& ontoMEDTRAD & OWL + iconique, Afrique de l'Ouest \\
\cline{2-3}
& ATHMed & Cadre proposé, Ghana \\
\cline{2-3}
& Deep Ontologies & OWL, Afrique \\
\cline{2-3}
& FLOPO, PO & OWL, botanique \\
\hline
Modèles de langage & gpt-6-luna ; Llama 3 8B & OpenWork (extraction) ; Ollama sur Colab et ngrok (RAG) \\
\hline
Plateforme & Google Colab & GPU T4 \\
\hline
Base de connaissances & Neo4j Desktop & Orientée graphe, Cypher \\
\hline
Orchestration RAG & LangChain & Consignes et appel des modèles \\
\hline
Interface web & Streamlit & Chatbot, recherche, fiches \\
\hline
Bibliothèques & Python & neo4j, langchain, streamlit, rdflib, etc. \\
\hline
\end{tabular}
\end{table}

% ------------------------------------------------------------
% II.1.5. Tableau de correspondance Matériel - Phases
% ------------------------------------------------------------

\subsection{Tableau de correspondance Matériel - Phases}
\label{subsec:correspondance_materiel}

Le tableau \ref{tab:correspondance_materiel} met en regard le matériel et les phases méthodologiques, et indique à quel moment et dans quel but chaque ressource entre en jeu.

\begin{table}[H]
\centering
\caption{Correspondance entre le matériel et les phases}
\label{tab:correspondance_materiel}
\begin{tabular}{|L{4cm}|L{3cm}|L{6cm}|}
\hline
\textbf{Matériel} & \textbf{Phase} & \textbf{Utilisation} \\
\hline
Corpus documentaire & Phase 3 & Constitution du corpus \\
\hline
Ontologies candidates & Phase 1 & Comparaison et sélection \\
\hline
gpt-6-luna & Phase 5 & Extraction \\
\hline
Llama 3 8B & Phase 7, Phase 8 & Génération RAG, évaluation \\
\hline
Google Colab & Phase 7, Phase 8 & Exécution de Llama 3 8B \\
\hline
Neo4j Desktop & Phase 6 & Base de connaissances \\
\hline
LangChain & Phase 7 & Orchestration RAG \\
\hline
Streamlit & Phase 7 & Interface web \\
\hline
Bibliothèques Python & Phase 5, 6, 7 & Développement \\
\hline
\end{tabular}
\end{table}

% ============================================================
% II.2. MÉTHODOLOGIE
% ============================================================

\section{Méthodologie}
\label{sec:methodologie}

Nous avons organisé la démarche en huit phases opérationnelles. Nous avons comparé les ontologies, défini les critères de recherche et constitué le corpus. Nous avons adapté l'ontologie, extrait les connaissances, construit la base de connaissances, développé le système d'accès et l'application web, puis évalué l'ensemble. Nous avons choisi l'ontologie avant de constituer le corpus afin de fixer le cadre conceptuel et de guider l'extraction.

% ------------------------------------------------------------
% II.2.1. Phase 1 : Identification et étude comparative des ontologies
% ------------------------------------------------------------

\subsection{Phase 1 : Identification et étude comparative des ontologies}
\label{subsec:phase1}

\textbf{Objectif} : Sélectionner l'ontologie la plus pertinente pour notre contexte.

Nous avons comparé les ontologies en trois temps. Nous avons défini une grille de critères, attribué une note à chacune, puis retenu celle qui a obtenu le meilleur score. Ce choix a fixé le modèle conceptuel des étapes suivantes.

\subsubsection{Définition de la grille de critères}

Nous avons retenu huit critères adaptés à nos besoins : couverture du domaine, richesse sémantique, disponibilité, documentation, interopérabilité, adaptabilité au contexte burkinabè, traçabilité et prise en compte des savoirs locaux. Le tableau \ref{tab:grille_criteres} présente ces critères.

\begin{table}[H]
\centering
\caption{Grille de critères de comparaison}
\label{tab:grille_criteres}
\begin{tabular}{|C{0.5cm}|L{3cm}|L{6.6cm}|C{2.5cm}|}
\hline
\textbf{N°} & \textbf{Critère} & \textbf{Description} & \textbf{Pondération} \\
\hline
1 & Couverture du domaine & Capacité à représenter plantes, usages, propriétés, maladies & Élevée (3) \\
\hline
2 & Richesse sémantique & Nombre et qualité des classes, relations, axiomes & Élevée (3) \\
\hline
3 & Disponibilité & Ontologie accessible et téléchargeable & Élevée (3) \\
\hline
4 & Documentation & Qualité de la documentation et des exemples & Moyenne (2) \\
\hline
5 & Interopérabilité & Compatibilité avec OWL, RDF, SPARQL & Moyenne (2) \\
\hline
6 & Adaptabilité & Facilité d'adaptation au contexte burkinabè & Élevée (3) \\
\hline
7 & Traçabilité & Capacité à associer les connaissances aux sources & Élevée (3) \\
\hline
8 & Prise en compte des savoirs locaux & Intégration des noms vernaculaires, recettes multi-plantes & Élevée (3) \\
\hline
\end{tabular}
\end{table}

Nous avons attribué à chaque critère une pondération proportionnelle à son importance. Les critères fortement pondérés (3) étaient déterminants pour représenter le contexte burkinabè, tandis que les critères moyennement pondérés (2) autorisaient davantage de compromis. Nous avons calculé le score global en multipliant chaque note par la pondération correspondante, en additionnant les produits et en divisant cette somme par la somme des pondérations. Le résultat est rapporté sur 5.

\subsubsection{Notation des ontologies}

Nous avons noté chaque ontologie de 1 (très faible) à 5 (très élevé) pour chacun des critères. Le tableau \ref{tab:notation_ontologies} présente les résultats.

\begin{table}[H]
\centering
\caption{Notation des ontologies candidates}
\label{tab:notation_ontologies}
\begin{tabular}{|L{3.5cm}|C{0.8cm}|C{0.8cm}|C{0.8cm}|C{0.8cm}|C{0.8cm}|C{0.8cm}|C{0.8cm}|C{0.8cm}|C{1.5cm}|}
\hline
\textbf{Ontologie} & \textbf{C1} & \textbf{C2} & \textbf{C3} & \textbf{C4} & \textbf{C5} & \textbf{C6} & \textbf{C7} & \textbf{C8} & \textbf{Score} \\
\hline
WATRIMed & 4 & 4 & 2 & 3 & 5 & 4 & 3 & 5 & 3.7 \\
\hline
ontoMEDTRAD & 5 & 5 & 5 & 4 & 4 & 5 & 4 & 5 & 4.7 \\
\hline
ATHMed & 4 & 4 & 3 & 3 & 4 & 3 & 2 & 3 & 3.2 \\
\hline
Deep Ontologies & 3 & 4 & 2 & 2 & 3 & 3 & 2 & 4 & 2.9 \\
\hline
FLOPO / PO & 2 & 3 & 5 & 4 & 5 & 2 & 1 & 1 & 2.7 \\
\hline
\end{tabular}
\end{table}

Nous avons analysé les notes afin de comparer la couverture du domaine et la prise en compte des savoirs locaux par les ontologies africaines. Nous avons également examiné la disponibilité et l'interopérabilité des ontologies botaniques, ainsi que les différences entre WATRIMed et ontoMEDTRAD concernant la richesse du modèle et l'accessibilité de la documentation.

Nous avons accordé une attention particulière à la disponibilité (C3) et à la documentation (C4). Nous avons évalué les ressources effectivement accessibles. WATRIMed est décrite par deux articles de congrès, dont l'un compte deux pages, mais son fichier d'ontologie n'est pas téléchargeable. La thèse consacrée à ontoMEDTRAD fournit la liste des classes et des propriétés, les métriques ainsi que l'origine des concepts.

\subsubsection{Sélection de l'ontologie}

WATRIMed \cite{watrimed2019, watrimed_design_2019} pourrait se distinguer par la richesse de son modèle documenté, qui compte 472 concepts et 75 propriétés, ainsi que par son alignement sur l'ontologie fondationnelle BioTopLite2. Toutefois, les deux publications n'en décrivent que les classes supérieures et les relations principales, et le fichier de l'ontologie n'est pas disponible pour une réutilisation directe.

Nous avons retenu ontoMEDTRAD \cite{kouame_systeme_2018}, qui a obtenu le meilleur score. Cette ontologie couvre l'Afrique de l'Ouest, représente les recettes multi-plantes, leurs ingrédients, les parties utilisées et les modes de préparation, et documente l'origine de ses concepts à partir de quinze ressources terminologiques et ontologiques. Sa thèse est accessible. Nous avons aussi pris en compte ses limites : sa validation globale reste annoncée comme perspective et son jeu de données comprend 22 plantes et 31 remèdes antipaludéens \cite{kouame_visual_2020}.

L'ontologie retenue fournira le cadre conceptuel pour extraire et organiser les connaissances. Elle guidera la constitution du corpus et son adaptation au contexte burkinabè.

\textbf{Livrable} : Ontologie ontoMEDTRAD sélectionnée.

% ------------------------------------------------------------
% II.2.2. Phase 2 : Définition des critères de recherche et de sélection
% ------------------------------------------------------------

\subsection{Phase 2 : Définition des critères de recherche et de sélection des documents}
\label{subsec:phase2}

\textbf{Objectif} : Définir les critères qui guideront la recherche et la sélection des documents scientifiques.

Nous avons défini deux catégories de critères, l'une pour rechercher les informations et l'autre pour sélectionner les documents. Elles visaient à assurer la pertinence, la qualité et la cohérence des sources retenues.

\subsubsection{Critères de recherche}

Nous avons retenu les critères suivants pour guider la recherche :

\begin{itemize}
\item \textbf{Les plantes} : espèces végétales documentées au Burkina Faso, avec identification botanique (nom scientifique, famille, synonymes).
\item \textbf{Les parties utilisées} : feuilles, racines, écorces, fruits, graines, fleurs, tiges, etc.
\item \textbf{Les pathologies et usages médicinaux} : paludisme, infections respiratoires, digestives, diabète, hypertension, etc.
\item \textbf{Les modes de préparation} : décoction, macération, infusion, poudre, extrait alcoolique, etc.
\item \textbf{Les niveaux de preuve} : savoirs traditionnels, études in vitro, études in vivo, essais cliniques.
\item \textbf{Les propriétés pharmacologiques} : antimicrobiennes, antipaludiques, anti-inflammatoires, antioxydantes, toxicité, etc.
\item \textbf{La composition chimique} : alcaloïdes, flavonoïdes, tanins, saponines, huiles essentielles, etc.
\end{itemize}

Ils couvrent les principales dimensions des connaissances sur les plantes médicinales et ont permis de cibler les informations utiles à l'étude.

\subsubsection{Critères de sélection des documents}

Nous avons sélectionné les documents selon les critères suivants :

\begin{itemize}
\item \textbf{Pertinence} : le document traite des plantes médicinales du Burkina Faso.
\item \textbf{Identification de la plante} : le nom scientifique est précisé.
\item \textbf{Disponibilité d'une référence scientifique} : le document est publié dans une revue à comité de lecture, une thèse, un mémoire ou un rapport de recherche.
\item \textbf{Fiabilité} : les sources sont vérifiables et citées.
\item \textbf{Actualité} : les documents récents sont privilégiés, sans exclure les travaux fondateurs.
\end{itemize}

Ces critères nous ont permis d'écarter les documents non scientifiques ou non pertinents et de constituer un corpus cohérent pour l'étude.

\textbf{Livrable} : grille de critères de recherche et de sélection.

% ------------------------------------------------------------
% II.2.3. Phase 3 : Constitution du corpus scientifique
% ------------------------------------------------------------

\subsection{Phase 3 : Constitution du corpus scientifique}
\label{subsec:phase3}

\textbf{Objectif} : Rassembler un corpus documentaire fiable et pertinent sur les plantes médicinales du Burkina Faso.

Nous avons constitué le corpus en quatre étapes. Nous avons recherché les documents dans Google Scholar, PubMed et ResearchGate, puis consulté burkinaherbiermed.org et les données de l'INERA. Nous avons inclus les documents portant sur les plantes médicinales du Burkina Faso, précisant le nom scientifique des plantes et publiés ou validés par une institution scientifique avec des sources vérifiables. Nous avons écarté les documents non scientifiques, ceux sans identification botanique et les documents anciens sans valeur de référence. Nous avons classé les documents par type, plante, pathologie et région et les avons enregistrés dans un répertoire structuré. Le corpus comprend 35 documents scientifiques et les données de burkinaherbiermed.org et de l'INERA.

\textbf{Livrable} : un corpus de 35 documents scientifiques, complété par les données de burkinaherbiermed.org et de l'INERA.

% ------------------------------------------------------------
% II.2.4. Phase 4 : Adaptation de l'ontologie sélectionnée
% ------------------------------------------------------------

\subsection{Phase 4 : Adaptation de l'ontologie sélectionnée}
\label{subsec:phase4}

\textbf{Objectif} : Adapter le modèle conceptuel de l'ontologie sélectionnée au contexte burkinabè.

Nous avons examiné les concepts existants, ajouté des classes et des propriétés, modifié certaines relations et documenté les changements.

\subsubsection{Analyse des concepts existants}

Nous avons étudié la structure de l'ontologie ontoMEDTRAD telle qu'elle est décrite dans la thèse \cite{kouame_systeme_2018}. L'ontologie compte 144 concepts, 76 propriétés associatives, 16 propriétés de données, 789 individus et 2 734 axiomes. Elle est formalisée en logique de description ALCHOIQ, éditée sous Protégé~5 et vérifiée par les raisonneurs HermiT et Pellet.

Le fichier OWL de l'ontologie n'étant pas téléchargeable publiquement, nous avons reconstitué sa structure à partir de la thèse. Nous nous sommes appuyés sur les captures d'écran de Protégé (figures 5.15 à 5.20) et les tableaux récapitulatifs des concepts, propriétés et métriques (tableaux 5.5, 5.21 et 5.22). Cette reconstitution a servi de base à l'adaptation présentée au chapitre \ref{chap:conception_developpement}.

Nous avons reconstitué la structure conceptuelle de l'ontologie, ses classes, propriétés et relations, mais pas ses instances d'origine. Nous n'avons pas réutilisé les 789 individus d'ontoMEDTRAD, liés au contexte et aux sources des travaux d'origine. Cette phase a produit le modèle conceptuel adapté ; le peuplement par les individus extraits de la littérature scientifique est intervenu en phase 5.

Ses concepts principaux décrivent le parcours de soin, du praticien et du patient aux symptômes, aux maladies, aux ressources médicinales, à l'examen clinique, aux recettes, aux ingrédients, aux parties utilisées, aux préparations et aux remèdes. L'ontologie distingue cinq types de plantes et neuf parties utilisées. Elle comprend aussi une classe pour les références scientifiques associées aux ressources, aux recettes et aux remèdes. Les concepts secondaires couvrent les contre-indications, les propriétés thérapeutiques, la toxicité, l'étiologie et la classification des maladies selon la CIM-10. Les relations de base expriment la subsomption et la relation partie-tout. Soixante-seize propriétés associatives complètent le modèle, dont la hiérarchie \textit{aPourForme}, utilisée pour l'iconisation des plantes.

Nous avons ensuite évalué ce modèle au regard du contexte burkinabè. Il couvre les concepts de base et les raisonneurs testés n'y ont signalé aucune incohérence, mais il nécessitait des adaptations pour représenter les noms vernaculaires en langues locales, les maladies spécifiques, les niveaux de preuve, les dosages et la traçabilité détaillée des sources. Sa validation globale reste une perspective, et son jeu de données se limite à 22 plantes et 31 remèdes antipaludéens.

\subsubsection{Ajout de classes}

La comparaison des fichiers du modèle conceptuel adapté et du schéma conservé dénombre 40 classes ajoutées. Le modèle adapté déclare également 25 propriétés objet au total, dont 20 absentes du schéma conservé, et 20 propriétés de données, dont 15 ajoutées. Les classes ajoutées couvrent notamment les noms vernaculaires, les usages thérapeutiques, les niveaux de preuve, les sources, les préparations, des parties végétales complémentaires, les composés chimiques, les activités pharmacologiques et les effets toxicologiques. Ces ajouts définissent le modèle conceptuel ; les individus et données extraits sont intégrés lors de la phase suivante.

Ces ajouts structurent des connaissances locales et scientifiques à différents niveaux. Ils ne préchargent pas les valeurs, comme les noms, pathologies, régions ou composés ; celles-ci proviennent des sources du corpus et sont enregistrées dans la base avec leur provenance.

\subsubsection{Ajout de propriétés}

Nous avons ajouté des propriétés de données pour les noms locaux, le niveau de preuve, le type d'étude, les références, le dosage et la toxicité, en complément de propriétés textuelles pour les noms vernaculaires et les descriptions d'activités ou d'effets. Des propriétés objet relient également l'usage à sa maladie cible, à la partie végétale utilisée, à la préparation et à la source ; d'autres décrivent les associations avec des composés, activités pharmacologiques, risques et localités. Les noms techniques exacts sont ceux déclarés dans la version courante du fichier Turtle.

L'enrichissement ainsi obtenu saisit des informations fines (noms vernaculaires, coordonnées, niveaux de preuve, références précises) que l'ontologie d'origine ne documentait pas pour le Burkina Faso.

\subsubsection{Modification de relations}

Nous avons structuré les liens entre plante, usage, maladie, partie utilisée et préparation, puis généralisé la traçabilité documentaire. La propriété objet \textit{estDocumenteePar} relie les ressources du graphe à une source scientifique; des sous-propriétés spécialisées permettent de préciser si la source documente une plante, un usage, un nom vernaculaire, une activité, un composé ou un effet toxicologique. \textit{aPourNiveauDePreuve} rattache un usage au niveau déclaré par la source. Ces liens ne signifient pas que les usages ou activités constituent des recommandations thérapeutiques ni une validation clinique.

Ces modifications assurent la traçabilité des connaissances et permettent d'exprimer des relations que le modèle d'origine, fondé sur la relation « est un » et sur la relation partie-tout, ne couvrait pas.

\subsubsection{Documentation des adaptations}

Nous avons consigné toutes les modifications dans un journal : chaque ajout est justifié, chaque choix est explicité, et les classes, propriétés et relations ajoutées, modifiées ou supprimées y sont décrites une à une.

Cette documentation n'est pas une formalité. Elle rend la démarche reproductible par un tiers et facilite les évolutions ultérieures de l'ontologie.

\textbf{Livrable} : Le modèle conceptuel adapté de l'ontologie, sans les données instanciées.

% ------------------------------------------------------------
% II.2.5. Phase 5 : Enrichissement de l'ontologie
% ------------------------------------------------------------

\subsection{Phase 5 : Enrichissement des connaissances à partir de la littérature scientifique}
\label{subsec:phase5}

\textbf{Objectif} : Extraire les connaissances du corpus pour alimenter la base organisée selon l'ontologie.

À partir du modèle conceptuel adapté en phase 4, nous avons enrichi l'ontologie avec des individus et des données extraits des sources scientifiques. Nous avons préparé ces données structurées et sourcées pour leur exploitation dans la base orientée graphe Neo4j.

Nous avons conçu les consignes d'extraction, procédé à l'extraction assistée par modèle de langue, contrôlé les données obtenues, puis structuré et préparé ces données pour leur intégration dans la base.

\subsubsection{Conception des prompts d'extraction}

Nos consignes d'extraction ciblaient les noms scientifiques et vernaculaires, la famille botanique, les parties utilisées, les usages, les préparations, les activités pharmacologiques, le niveau de preuve et la source. Le chapitre \ref{chap:conception_developpement} présente un exemple de consigne, avec le format JSON et les champs demandés.

Nous avons employé une même consigne pour les documents traités. Les champs absents reçoivent la valeur \texttt{null} ou une liste vide, et aucune valeur n'est complétée par déduction.

\subsubsection{Exécution de l'extraction}

Nous avons soumis les documents à gpt-6-luna dans l'environnement OpenWork, un par un, et enregistré les résultats en JSON. L'extraction a porté sur dix des 35 documents du corpus, principalement des enquêtes ethnobotaniques ; les 25 autres n'ont pas été extraits.

\subsubsection{Contrôle de la qualité des extractions}

Nous avons repris le niveau de preuve indiqué par la source lorsqu'il était explicite, qu'il s'agisse d'un savoir traditionnel, d'une étude \textit{in vitro}, d'une étude \textit{in vivo} ou d'un essai clinique. Nous n'avons pas évalué indépendamment la solidité des faits et avons laissé le champ vide lorsque la source ne précisait pas ce niveau. Un expert externe n'a pas relu les extractions. Nous avons contrôlé leur structure, puis comparé une partie des faits aux documents d'origine à la phase 8.

\subsubsection{Structuration des données}

Nous avons organisé les données selon le modèle de l'ontologie dans un classeur d'individus, par type de connaissance (sources, plantes, noms vernaculaires, usages, préparations, maladies, activités, composés, toxicologie et localités), puis les avons transformées en individus au format Turtle.

\subsubsection{Préparation de l'intégration dans la base de connaissances}

Le fichier OntoPlantMedBF-peuplee.ttl contient des individus extraits de dix documents scientifiques, en plus du schéma adapté. Ce fichier constitue l'ontologie enrichie importée dans Neo4j à la phase 6. Les données de burkinaherbiermed.org et de l'INERA, consultées lors de la constitution du corpus, n'y ont pas été intégrées.

\textbf{Livrable} : OntoPlantMedBF, ontologie enrichie avec les données extraites et leurs sources.

% ------------------------------------------------------------
% II.2.6. Phase 6 : Construction et interrogation de la base
% ------------------------------------------------------------

\subsection{Phase 6 : Construction et interrogation de la base de connaissances}
\label{subsec:phase6}

\textbf{Objectif} : Importer et exploiter dans Neo4j l'ontologie réajustée \textit{OntoPlantMedBF}.

Nous avons importé dans Neo4j l'ontologie réajustée \textit{OntoPlantMedBF} et ses données, en les représentant sous forme de nœuds, propriétés et relations du graphe. Nous avons vérifié la correspondance avec le modèle ontologique et testé les requêtes de la base.

\subsubsection{Import et ajustement du modèle dans le graphe}

Nous avons importé les classes et propriétés d'OntoPlantMedBF dans Neo4j sous forme de labels, propriétés et relations. Les principaux nœuds représentent les plantes, maladies, usages, parties, préparations, sources, composés chimiques et régions. Les relations du graphe reprennent les liens du modèle, avec des noms adaptés aux conventions Neo4j, et les propriétés conservent notamment les noms, niveaux de preuve, DOI et coordonnées géographiques.

Nous avons ainsi importé dans le graphe la structure réajustée d'OntoPlantMedBF et ses données. Les entités sont représentées par des nœuds, leurs liens par des relations et leurs caractéristiques par des propriétés, conformément au modèle ontologique.

\subsubsection{Import des données}

Nous avons chargé les données préparées, créé ou rapproché les nœuds et relations au moyen de requêtes Cypher, puis vérifié l'intégrité et la cohérence de l'import avec le modèle ontologique.

Après l'import, nous avons contrôlé que les données, les relations et les propriétés avaient été correctement reprises.

\subsubsection{Requêtes de test}

Nous avons testé quatre familles de requêtes : les requêtes simples sur les plantes associées à une maladie, les requêtes complexes sur les recettes multi-plantes, les requêtes de traçabilité et les requêtes spatiales par région.

Ces vérifications nous ont permis de contrôler le fonctionnement de la base et les résultats renvoyés par les requêtes.

\textbf{Livrable} : Une base Neo4j peuplée et interrogeable, complémentaire d'OntoPlantMedBF.

% ------------------------------------------------------------
% II.2.7. Phase 7 : Exploitation par l'IA et développement web
% ------------------------------------------------------------

\subsection{Phase 7 : Exploitation par IA et développement de l'application web}
\label{subsec:phase7}

\textbf{Objectif} : Développer le système d'accès intelligent et l'interface utilisateur.

Nous avons mis en place le système RAG en définissant son architecture, en construisant la chaîne de traitement et en intégrant ses composants.

\subsubsection{Architecture du système}

L'architecture relie l'utilisateur à Streamlit, LangChain, Neo4j et Llama 3 jusqu'à la réponse. Elle associe la recherche dans la base à la génération par le modèle de langue et conserve le lien entre les informations et leurs sources. Le chapitre \ref{chap:conception_developpement} décrit la chaîne de traitement, de l'interprétation de la question à l'affichage des références.

\subsubsection{Intégration des composants}

Nous avons connecté l'application à Neo4j et au modèle de langue, puis configuré les consignes de génération des réponses. Cette intégration a réuni les composants du système RAG.

\subsubsection{Développement de l'application web}

Nous avons conçu l'application autour d'une page d'accueil avec chatbot, d'une recherche par plante, d'une recherche par pathologie et de fiches descriptives. Nous y avons intégré le système RAG et la gestion de l'historique. L'interface permet d'interroger la base en langage naturel et de recevoir des réponses reliées à leurs sources. Le chapitre \ref{chap:conception_developpement} en décrit la conception et les composants.

\subsubsection{Tests et validation}

Nous avons mené des tests fonctionnels pour vérifier le fonctionnement de l'application et des tests utilisateurs pour examiner sa facilité d'utilisation et la pertinence des réponses.

\textbf{Livrable} : Un système RAG fonctionnel et une application web fonctionnelle.

% ------------------------------------------------------------
% II.2.8. Phase 8 : Évaluation et validation
% ------------------------------------------------------------

\subsection{Phase 8 : Évaluation et validation}
\label{subsec:phase8}

\textbf{Objectif} : Évaluer la qualité de l'ontologie, du système RAG et de l'application web.

Nous avons évalué trois composantes : l'ontologie, le système RAG et l'application web.

\subsubsection{Évaluation de l'ontologie}

Nous avons évalué l'ontologie selon les quatre critères présentés dans le tableau \ref{tab:eval_ontologie}.

\begin{table}[H]
\centering
\caption{Critères d'évaluation de l'ontologie}
\label{tab:eval_ontologie}
\begin{tabular}{|L{3cm}|L{6cm}|L{4cm}|}
\hline
\textbf{Critère} & \textbf{Description} & \textbf{Méthode} \\
\hline
Couverture & Capacité à représenter les connaissances du corpus & Vérification manuelle \\
\hline
Cohérence & Absence de contradictions & Requêtes de test \\
\hline
Traçabilité & Association des connaissances aux sources & Vérification des relations \\
\hline
Extensibilité & Facilité d'ajout de nouvelles connaissances & Test d'ajout \\
\hline
\end{tabular}
\end{table}

\subsubsection{Évaluation du système RAG}

Nous avons évalué le système RAG selon les quatre critères présentés dans le tableau \ref{tab:eval_rag}.

\begin{table}[H]
\centering
\caption{Critères d'évaluation du système RAG}
\label{tab:eval_rag}
\begin{tabular}{|L{3cm}|L{6cm}|L{4cm}|}
\hline
\textbf{Critère} & \textbf{Description} & \textbf{Métrique} \\
\hline
Pertinence & Les réponses répondent à la question & Précision \\
\hline
Fidélité & Les réponses sont fidèles aux sources & Rappel \\
\hline
Exactitude & Les réponses sont correctes & F1-score \\
\hline
Traçabilité & Les sources sont affichées & Taux de couverture \\
\hline
\end{tabular}
\end{table}

Nous avons constitué un jeu de 50 questions de référence portant sur les plantes, les usages et leurs sources, puis comparé les réponses du système aux réponses attendues.

\subsubsection{Évaluation de l'application web}

Nous avons évalué l'application web selon les quatre critères présentés dans le tableau \ref{tab:eval_web}.

\begin{table}[H]
\centering
\caption{Critères d'évaluation de l'application web}
\label{tab:eval_web}
\begin{tabular}{|L{3cm}|L{6cm}|L{4cm}|}
\hline
\textbf{Critère} & \textbf{Description} & \textbf{Méthode} \\
\hline
Facilité d'utilisation & L'interface est intuitive & Questionnaire \\
\hline
Compréhension & Les réponses sont claires & Questionnaire \\
\hline
Utilité & L'application répond aux besoins & Questionnaire \\
\hline
Performance & Temps de réponse & Mesure \\
\hline
\end{tabular}
\end{table}

Nous avons sollicité des étudiants, des chercheurs et des tradipraticiens pour cette évaluation.

\textbf{Livrable} : Rapport d'évaluation et de validation.

% ============================================================
% II.3. SYNTHÈSE ET TRANSITION
% ============================================================

\section{Synthèse et transition}
\label{sec:synthese}

\subsection{Tableau de correspondance Phases - Livrables}

Le tableau \ref{tab:correspondance_phases} récapitule les objectifs et les livrables des huit phases.

\begin{table}[H]
\centering
\caption{Correspondance entre les phases, les objectifs et les livrables}
\label{tab:correspondance_phases}
\begin{tabular}{|L{2.2cm}|L{4.5cm}|L{6.8cm}|}
\hline
\textbf{Phase} & \textbf{Objectif} & \textbf{Livrable} \\
\hline
Phase 1 & Sélectionner l'ontologie la plus pertinente & Ontologie ontoMEDTRAD sélectionnée \\
\hline
Phase 2 & Définir les critères de recherche et de sélection & Grille de critères \\
\hline
Phase 3 & Rassembler un corpus fiable et pertinent & Corpus de 35 documents \\
\hline
Phase 4 & Adapter le modèle conceptuel de l'ontologie & Modèle conceptuel adapté \\
\hline
Phase 5 & Extraire les connaissances et enrichir l'ontologie & Ontologie enrichie en données \\
\hline
Phase 6 & Importer OntoPlantMedBF dans la base graphe & OntoPlantMedBF intégrée à Neo4j, base interrogeable \\
\hline
Phase 7 & Développer le système RAG et l'application web & Système RAG et application web \\
\hline
Phase 8 & Évaluer la qualité du système & Rapport d'évaluation \\
\hline
\end{tabular}
\end{table}

% ============================================================
% CONCLUSION DU CHAPITRE
% ============================================================

\section*{Conclusion du chapitre}
\addcontentsline{toc}{section}{Conclusion du chapitre}

Ce chapitre a exposé le matériel et la méthodologie de notre recherche. Le matériel réunit un corpus de 35 documents scientifiques, les données de burkinaherbiermed.org et de l'INERA, cinq ontologies candidates et un environnement technique bâti sur gpt-6-luna pour l'extraction, Llama 3 8B sur Google Colab pour le RAG, Neo4j Desktop, LangChain et Streamlit.

Nous avons suivi huit phases : comparaison des ontologies, définition des critères de recherche et de sélection, constitution du corpus, adaptation de l'ontologie, extraction des connaissances, construction et interrogation de la base, développement du système d'accès et de l'application, puis évaluation. Nous avons décrit les actions menées et les livrables associés.

Nous avons choisi l'ontologie avant de constituer le corpus afin d'assurer la cohérence entre le modèle conceptuel et les données collectées, et de cibler les informations pertinentes dès la collecte.

Nous signalons deux limites. Les niveaux de preuve associés aux connaissances proviennent des sources et n'ont pas fait l'objet d'une évaluation indépendante. De plus, un expert externe n'a pas validé les extractions. Ces limites sont à prendre en compte dans les travaux ultérieurs. Nous avons néanmoins conservé le lien entre chaque information et sa source.

Le chapitre \ref{chap:resultats_discussion} présente les résultats de l'évaluation menée en phase 8.

Le chapitre suivant présente la conception et le développement du système, notamment l'ontologie adaptée, la base de connaissances, le système RAG et l'application web.
).

\phantomsection
\section*{Introduction du chapitre}

Après avoir posé le cadre conceptuel et exposé l'état de l'art, nous présentons le matériel mobilisé et la démarche méthodologique suivie. Ce chapitre fait le lien entre la réflexion théorique et la mise en œuvre technique.

Nous présentons le matériel mobilisé, notamment le corpus documentaire, les ontologies candidates et l'environnement technique. Nous décrivons les huit phases de la démarche suivie, de l'étude comparative des ontologies à l'évaluation du système, puis en proposons une synthèse.

% ============================================================
% II.1. MATÉRIEL
% ============================================================

\section{Matériel}
\label{sec:materiel}

Le matériel comprend le corpus documentaire, les ontologies candidates et l'environnement technique. Les deux tableaux qui suivent récapitulent ces ressources, réparties entre documentation, modèles conceptuels et outils techniques.

% ------------------------------------------------------------
% II.1.1. Corpus documentaire
% ------------------------------------------------------------

\subsection{Corpus documentaire}
\label{subsec:corpus}

Nous avons choisi l'ontologie avant de constituer le corpus afin de déterminer les notions et les relations qui guideraient l'extraction. Nous avons réuni les publications scientifiques et les données institutionnelles consultées sur burkinaherbiermed.org et auprès de l'INERA.

La figure \ref{fig:flux_corpus} résume le processus de constitution du corpus, du choix de l'ontologie à la sélection des documents.

\begin{figure}[H]
\centering
\includegraphics[width=0.8\textwidth]{flux_corpus.png}
\caption{Constitution du corpus documentaire}
\label{fig:flux_corpus}
\end{figure}

\subsubsection{Sources scientifiques publiées}

Notre corpus scientifique comprend 35 rapports de recherche, articles, thèses et mémoires provenant de revues, d'universités et d'institutions de recherche burkinabè. Ces travaux ethnobotaniques, ethnopharmacologiques et phytochimiques documentent les noms scientifiques et vernaculaires, les parties utilisées, les pathologies traitées, les modes de préparation et, pour les études qui les rapportent, les dosages et les données de toxicité.

Nous avons retenu ces travaux pour leur qualité et leur pertinence au contexte burkinabè, des publications fondatrices aux plus récentes. Cette diversité d'espèces, de pathologies et de pratiques thérapeutiques soutient l'extraction. Quatre travaux illustrent les types de données mobilisés dans le corpus :
\begin{itemize}
\item Zerbo et al.~\cite{zerbo_plantes_2011} ont étudié les plantes médicinales chez les Sanan du Burkina Faso. Leur travail documente 85 espèces utilisées pour traiter 63 affections.
\item Ouédraogo et al.~\cite{ouedraogo_revue_2024} ont dressé une revue des plantes médicinales à potentiel aromatique du pays, recensant 42 espèces et en caractérisant les propriétés et les usages.
\item Tibiri et al.~\cite{tibiri2020} ont mené une enquête à l'échelle nationale sur les plantes utilisées contre les maladies hépatiques, documentant plus de 2 000 entrées de plantes correspondant à 212 espèces.
\item Ouédraogo et al.~\cite{ouedraogo2020markets} ont enquêté auprès de 30 commerçants de plantes médicinales dans six grandes villes, identifiant 104 espèces utilisées pour diverses pathologies.
\end{itemize}

Ces travaux fournissent des méthodes réutilisables et des observations de terrain recueillies au Burkina Faso.

\subsubsection{Données de l'Herbier National des Plantes Médicinales du Burkina Faso}

Le site burkinaherbiermed.org diffuse des données sur les plantes médicinales du pays, notamment des informations botaniques, des usages thérapeutiques et des références bibliographiques. Ces données proviennent du projet « Herbier avec chimiothèque », piloté par le Ministère de la Santé. Ce projet a pour objectif de constituer un herbier national dédié aux plantes médicinales et une chimiothèque de marqueurs chimiques. Il documente les plantes utilisées contre les dix pathologies prioritaires définies par le ministère, géolocalise les stations et conserve des spécimens. Ces informations alimentent une base nationale accessible par application mobile et site web. Le projet associe un herbier, une chimiothèque et une base de données ethnobotaniques et géographiques. Il complète les autres herbiers burkinabè, notamment le HNBU du CNRST/INERA, le Centre National de Semences Forestières et l'herbier de l'Université Joseph Ki-Zerbo.

\subsubsection{Données de l'INERA}

La base de données de l'Institut de l'Environnement et de Recherches Agricoles (INERA) comprend plus de 7 000 espèces végétales et les coordonnées géographiques des stations où elles ont été observées. Ces données complètent les informations ethnobotaniques par une dimension spatiale et permettent d'étudier la répartition des plantes au Burkina Faso. Le rapprochement des usages rapportés et de la distribution géographique peut faire apparaître des régularités régionales, qui restent à vérifier.

Le tableau \ref{tab:corpus} récapitule les sources du corpus.

\begin{table}[H]
\centering
\caption{Sources du corpus documentaire}
\label{tab:corpus}
\begin{tabular}{|L{4cm}|L{2.8cm}|L{3cm}|L{5cm}|}
\hline
\textbf{Source} & \textbf{Type} & \textbf{Volume estimé} & \textbf{Contenu principal} \\
\hline
Articles et thèses & Publications scientifiques & 35 documents & Études ethnobotaniques, pharmacologiques \\
\hline
burkinaherbiermed.org & Base de données en ligne & Plusieurs centaines de plantes & Données botaniques, usages, références \\
\hline
INERA & Base de données institutionnelle & > 7 000 espèces & Coordonnées géographiques, identification \\
\hline
\end{tabular}
\end{table}

Les documents du corpus relèvent de types d'étude différents, et chacun est exploitable pour les informations qu'il rapporte. Trois profils se distinguent. Les enquêtes ethnobotaniques et ethnopharmacologiques, menées auprès de tradipraticiens ou dans les marchés, rapportent les noms vernaculaires, les parties utilisées, les modes de préparation et les pathologies traitées, souvent plante par plante. Les études expérimentales, notamment les essais \textit{in vitro} et les travaux phytochimiques, mesurent des activités pharmacologiques, des concentrations, des niveaux de preuve et parfois les lieux de récolte. Elles renseignent moins sur les noms vernaculaires et les pratiques de préparation. Enfin, les revues de synthèse dressent des inventaires à grande échelle et servent surtout à repérer les sources primaires.

Nous avons appliqué trois règles à l'extraction. Chaque source ne renseignait que les champs qu'elle mesurait. Nous avons laissé vides les champs non couverts, sans les compléter par déduction. Nous n'avons pas retenu comme fait établi un usage traditionnel rapporté dans un article expérimental sans référence ethnobotanique. Enfin, nous n'avons pas additionné les nombres d'espèces d'une source à l'autre, car une revue peut reprendre les espèces documentées dans les enquêtes qu'elle cite. Le tableau \ref{tab:corpus} présente les sources document par document, sans agréger les volumes.

Ce corpus rassemble des études menées dans plusieurs régions du Burkina Faso, portant sur différentes pathologies et mobilisant des approches méthodologiques variées. Il a servi de base à l'extraction des connaissances et à l'enrichissement de l'ontologie. Le choix préalable de l'ontologie a guidé l'extraction et l'organisation des documents selon les concepts et les relations définis.

% ------------------------------------------------------------
% II.1.2. Ontologies candidates
% ------------------------------------------------------------

\subsection{Ontologies candidates}
\label{subsec:ontologies_candidates}

Nous avons retenu cinq ontologies candidates pour l'étude comparative. Le chapitre \ref{chap:cadre_theorique}, aux sections \ref{subsubsec:ontologies_internationales} et \ref{subsubsec:ontologies_africaines}, les présente. Elles couvrent différents domaines, notamment la botanique, la médecine traditionnelle, la pharmacologie et l'ethnomédecine. Le tableau \ref{tab:ontologies_candidates} en présente les caractéristiques.

Nous les avons comparées à l'aide de la grille de critères présentée en phase 1 (section \ref{subsec:phase1}). Nous avons sélectionné l'ontologie la plus adaptée à notre contexte avant de constituer le corpus, afin de fixer le modèle conceptuel qui a guidé l'extraction.

\begin{table}[H]
\centering
\caption{Ontologies candidates pour la comparaison}
\label{tab:ontologies_candidates}
\begin{tabular}{|C{0.5cm}|L{3cm}|L{3cm}|L{3cm}|L{2cm}|L{2.5cm}|}
\hline
\textbf{N°} & \textbf{Ontologie} & \textbf{Domaine} & \textbf{Zone géographique} & \textbf{Format} & \textbf{Source d'accès} \\
\hline
1 & WATRIMed & Médecine traditionnelle & Afrique de l'Ouest & OWL, graphe de connaissances & IOS Press (2 articles) \\
\hline
2 & ontoMEDTRAD & Médecine traditionnelle & Afrique de l'Ouest & OWL + iconique & Thèse A. Kouamé (UGB) \\
\hline
3 & ATHMed & Médecine traditionnelle & Afrique (Ghana) & Cadre conceptuel & Frontiers in AI \\
\hline
4 & Deep Ontologies & Médecine traditionnelle & Afrique & OWL (Protégé) & Articles ARIMA \\
\hline
5 & FLOPO, PO & Botanique & International & OWL & OBO Foundry \\
\hline
\end{tabular}
\end{table}

Chacune présente des acquis et des lacunes. Les ontologies botaniques (FLOPO, PO) décrivent bien la plante mais ne traitent pas des usages médicinaux. Les ontologies pharmaceutiques (DrOn) se resserrent sur les médicaments modernes. L'OTM, prometteuse, reste sans ancrage africain. Les initiatives africaines (WATRIMed, ontoMEDTRAD, ATHMed) s'intéressent en revanche à la médecine traditionnelle du terrain, avec ses dimensions culturelles et ses recettes multi-plantes, même si ATHMed en reste, elle, au stade d'un cadre proposé. La phase 1 en livre la comparaison détaillée.

% ------------------------------------------------------------
% II.1.3. Outils et environnement technique
% ------------------------------------------------------------

\subsection{Outils et environnement technique}
\label{subsec:environnement}

Les outils techniques occupent une place importante dans notre démarche. Nous décrivons les phases concernées et les raisons de leur choix.

La figure \ref{fig:dispositif_technique} résume l'environnement d'expérimentation, en précisant où chaque composant est exécuté et comment ils communiquent.

\begin{figure}[H]
\centering
\includegraphics[width=0.8\textwidth]{dispositif_technique.png}
\caption{Dispositif technique d'expérimentation}
\label{fig:dispositif_technique}
\end{figure}

\subsubsection{Modèles de langage}

Nous avons utilisé deux modèles, chacun pour une étape distincte. Pour l'extraction des connaissances (phase 5), nous avons utilisé \textbf{gpt-6-luna} d'OpenAI, accessible par l'environnement OpenWork. Pour la génération des réponses du système RAG (phase 7), nous avons utilisé \textbf{Llama 3 8B} \cite{meta2024llama3}, un modèle ouvert. Gemini (\textit{gemini-3.5-flash-lite}) nous a servi de point de comparaison lors de l'évaluation.

Le tableau \ref{tab:llama_acces} présente les modes d'accès à ces modèles.

\begin{table}[H]
\centering
\caption{Modèles de langage et modes d'accès}
\label{tab:llama_acces}
\begin{tabular}{|L{4cm}|L{4cm}|L{6cm}|}
\hline
\textbf{Usage} & \textbf{Modèle} & \textbf{Mode d'accès} \\
\hline
Extraction (phase 5) & gpt-6-luna & OpenWork \\
\hline
Génération RAG (phase 7) & Llama 3 8B & Ollama sur Google Colab, API exposée par ngrok \\
\hline
Comparaison (phase 8) & gemini-3.5-flash-lite & API de Google \\
\hline
\end{tabular}
\end{table}

Nous avons retenu Llama 3 pour la génération parce qu'il s'agit d'un modèle ouvert que nous pouvons exécuter sur une instance que nous contrôlons. La version 8B tient dans la mémoire d'un GPU accessible gratuitement sur Colab. Ollama sert le modèle et ngrok expose son API à l'application web, qui reste sur le poste local.

\subsubsection{Plateforme d'exécution : Google Colab}

Nous avons exécuté Llama 3 8B sur \textbf{Google Colab}, une plateforme de notebooks donnant accès à des GPU adaptés à l'inférence.

\textbf{Phases d'utilisation :} phases 7 et 8 pour la génération des réponses et l'évaluation du système RAG.

Cette plateforme nous a donné accès à un GPU sans nécessiter l'acquisition d'une infrastructure locale. La session reste cependant temporaire et l'adresse exposée par ngrok change à chaque lancement, ce qui convient à un prototype mais pas à un service pérenne.

\subsubsection{Base de connaissances : Neo4j Desktop}

Nous avons utilisé \textbf{Neo4j Desktop}, la version locale de Neo4j, pour créer, gérer et interroger la base de connaissances orientée graphe.

\textbf{Phase d'utilisation :} phase 6 pour la construction et l'interrogation de la base.

Le tableau \ref{tab:neo4j} présente les caractéristiques de Neo4j Desktop.

\begin{table}[H]
\centering
\caption{Caractéristiques de Neo4j Desktop}
\label{tab:neo4j}
\begin{tabular}{|L{5cm}|L{8cm}|}
\hline
\textbf{Caractéristique} & \textbf{Valeur} \\
\hline
Type & Base de données orientée graphe \\
\hline
Langage de requête & Cypher \\
\hline
Version & Desktop (locale) \\
\hline
Format d'import & CSV, JSON \\
\hline
Interface & Neo4j Browser, Bloom \\
\hline
\end{tabular}
\end{table}

Nous avons choisi Neo4j pour représenter les relations entre les plantes, les usages, les propriétés et les sources, interroger le graphe avec Cypher et assurer la traçabilité des connaissances. La base permet également d'ajouter des connaissances et de visualiser les liens entre les entités.

\subsubsection{Orchestration RAG : LangChain}

Nous avons utilisé \textbf{LangChain} pour définir les consignes et appeler les modèles de langage de façon uniforme. La recherche dans Neo4j passe par le pilote officiel et par nos propres requêtes.

\textbf{Phase d'utilisation :} phase 7 pour l'accès aux connaissances et le développement web.

Le tableau \ref{tab:langchain} présente les composants LangChain utilisés.

\begin{table}[H]
\centering
\caption{Composants LangChain utilisés}
\label{tab:langchain}
\begin{tabular}{|L{5cm}|L{8cm}|}
\hline
\textbf{Composant} & \textbf{Rôle} \\
\hline
PromptTemplate & Préparation des consignes de génération \\
\hline
ChatOpenAI & Appel de Llama 3 par l'API d'Ollama, compatible avec celle d'OpenAI \\
\hline
ChatGoogleGenerativeAI & Appel de Gemini pour la comparaison \\
\hline
\end{tabular}
\end{table}

Nous avons retenu LangChain parce qu'il permet de changer de modèle sans modifier le reste du système. Nous n'avons pas confié la sélection des données à une chaîne générant du Cypher : le chapitre \ref{chap:conception_developpement} justifie ce choix \cite{lewis2020rag}.

\subsubsection{Interface web : Streamlit}

Nous avons développé l'application web avec \textbf{Streamlit}, un framework Python qui permet de créer des applications web interactives.

\textbf{Phase d'utilisation :} phase 7 pour le développement de l'application web.

L'application propose un chatbot en langage naturel, des recherches par plante, pathologie ou région, des fiches descriptives et l'affichage des sources et des niveaux de preuve. Nous avons également prévu une fonction de visualisation du graphe de connaissances.

Nous avons retenu Streamlit pour sa simplicité d'utilisation, sa gestion des sessions, son intégration avec LangChain et la rapidité de développement qu'il permet. Il facilite aussi le déploiement de l'application.

\subsubsection{Bibliothèques Python}

Nous avons mobilisé plusieurs bibliothèques Python au cours des phases 5, 6 et 7. Le tableau \ref{tab:bibliotheques} en présente la liste et le rôle.

\begin{table}[H]
\centering
\caption{Bibliothèques Python utilisées}
\label{tab:bibliotheques}
\begin{tabular}{|L{5cm}|L{3cm}|L{5cm}|}
\hline
\textbf{Bibliothèque} & \textbf{Version} & \textbf{Rôle} \\
\hline
\texttt{neo4j} & 5.28 & Pilote de connexion à Neo4j \\
\hline
\texttt{langchain}, \texttt{langchain-core} & 0.3 & Consignes et appel des modèles \\
\hline
\texttt{langchain-openai} & 0.3 & Accès à l'API d'Ollama \\
\hline
\texttt{langchain-google-genai} & 2.1 & Accès à Gemini \\
\hline
\texttt{streamlit} & 1.64 & Interface web \\
\hline
\texttt{plotly} & 7.1 & Visualisation du graphe \\
\hline
\texttt{rdflib} & 7.6 & Lecture du fichier Turtle \\
\hline
\texttt{pytest} & 9.1 & Tests logiciels \\
\hline
\texttt{pyngrok} & Sur Colab & Ouverture du tunnel ngrok \\
\hline
\end{tabular}
\end{table}

Ces bibliothèques couvrent les besoins de la recherche, notamment la connexion à la base, l'orchestration du système RAG, le développement de l'interface, le chargement du modèle, la manipulation de l'ontologie et le traitement des données.

% ------------------------------------------------------------
% II.1.4. Tableau récapitulatif du matériel
% ------------------------------------------------------------

\subsection{Tableau récapitulatif du matériel}
\label{subsec:recap_materiel}

Le tableau \ref{tab:recap_materiel} rassemble l'ensemble du matériel mobilisé.

\begin{table}[H]
\centering
\caption{Récapitulatif du matériel}
\label{tab:recap_materiel}
\begin{tabular}{|L{3cm}|L{4cm}|L{6cm}|}
\hline
\textbf{Catégorie} & \textbf{Élément} & \textbf{Spécification} \\
\hline
\multirow{3}{*}{Corpus} & Documents scientifiques & 35 rapports, articles, thèses \\
\cline{2-3}
& Base burkinaherbiermed.org & Plusieurs centaines de plantes \\
\cline{2-3}
& Base INERA & > 7 000 espèces géolocalisées \\
\hline
\multirow{5}{*}{Ontologies} & WATRIMed & OWL, Afrique de l'Ouest \\
\cline{2-3}
& ontoMEDTRAD & OWL + iconique, Afrique de l'Ouest \\
\cline{2-3}
& ATHMed & Cadre proposé, Ghana \\
\cline{2-3}
& Deep Ontologies & OWL, Afrique \\
\cline{2-3}
& FLOPO, PO & OWL, botanique \\
\hline
Modèles de langage & gpt-6-luna ; Llama 3 8B & OpenWork (extraction) ; Ollama sur Colab et ngrok (RAG) \\
\hline
Plateforme & Google Colab & GPU T4 \\
\hline
Base de connaissances & Neo4j Desktop & Orientée graphe, Cypher \\
\hline
Orchestration RAG & LangChain & Consignes et appel des modèles \\
\hline
Interface web & Streamlit & Chatbot, recherche, fiches \\
\hline
Bibliothèques & Python & neo4j, langchain, streamlit, rdflib, etc. \\
\hline
\end{tabular}
\end{table}

% ------------------------------------------------------------
% II.1.5. Tableau de correspondance Matériel - Phases
% ------------------------------------------------------------

\subsection{Tableau de correspondance Matériel - Phases}
\label{subsec:correspondance_materiel}

Le tableau \ref{tab:correspondance_materiel} met en regard le matériel et les phases méthodologiques, et indique à quel moment et dans quel but chaque ressource entre en jeu.

\begin{table}[H]
\centering
\caption{Correspondance entre le matériel et les phases}
\label{tab:correspondance_materiel}
\begin{tabular}{|L{4cm}|L{3cm}|L{6cm}|}
\hline
\textbf{Matériel} & \textbf{Phase} & \textbf{Utilisation} \\
\hline
Corpus documentaire & Phase 3 & Constitution du corpus \\
\hline
Ontologies candidates & Phase 1 & Comparaison et sélection \\
\hline
gpt-6-luna & Phase 5 & Extraction \\
\hline
Llama 3 8B & Phase 7, Phase 8 & Génération RAG, évaluation \\
\hline
Google Colab & Phase 7, Phase 8 & Exécution de Llama 3 8B \\
\hline
Neo4j Desktop & Phase 6 & Base de connaissances \\
\hline
LangChain & Phase 7 & Orchestration RAG \\
\hline
Streamlit & Phase 7 & Interface web \\
\hline
Bibliothèques Python & Phase 5, 6, 7 & Développement \\
\hline
\end{tabular}
\end{table}

% ============================================================
% II.2. MÉTHODOLOGIE
% ============================================================

\section{Méthodologie}
\label{sec:methodologie}

Nous avons organisé la démarche en huit phases opérationnelles. Nous avons comparé les ontologies, défini les critères de recherche et constitué le corpus. Nous avons adapté l'ontologie, extrait les connaissances, construit la base de connaissances, développé le système d'accès et l'application web, puis évalué l'ensemble. Nous avons choisi l'ontologie avant de constituer le corpus afin de fixer le cadre conceptuel et de guider l'extraction.

% ------------------------------------------------------------
% II.2.1. Phase 1 : Identification et étude comparative des ontologies
% ------------------------------------------------------------

\subsection{Phase 1 : Identification et étude comparative des ontologies}
\label{subsec:phase1}

\textbf{Objectif} : Sélectionner l'ontologie la plus pertinente pour notre contexte.

Nous avons comparé les ontologies en trois temps. Nous avons défini une grille de critères, attribué une note à chacune, puis retenu celle qui a obtenu le meilleur score. Ce choix a fixé le modèle conceptuel des étapes suivantes.

\subsubsection{Définition de la grille de critères}

Nous avons retenu huit critères adaptés à nos besoins : couverture du domaine, richesse sémantique, disponibilité, documentation, interopérabilité, adaptabilité au contexte burkinabè, traçabilité et prise en compte des savoirs locaux. Le tableau \ref{tab:grille_criteres} présente ces critères.

\begin{table}[H]
\centering
\caption{Grille de critères de comparaison}
\label{tab:grille_criteres}
\begin{tabular}{|C{0.5cm}|L{3cm}|L{6.6cm}|C{2.5cm}|}
\hline
\textbf{N°} & \textbf{Critère} & \textbf{Description} & \textbf{Pondération} \\
\hline
1 & Couverture du domaine & Capacité à représenter plantes, usages, propriétés, maladies & Élevée (3) \\
\hline
2 & Richesse sémantique & Nombre et qualité des classes, relations, axiomes & Élevée (3) \\
\hline
3 & Disponibilité & Ontologie accessible et téléchargeable & Élevée (3) \\
\hline
4 & Documentation & Qualité de la documentation et des exemples & Moyenne (2) \\
\hline
5 & Interopérabilité & Compatibilité avec OWL, RDF, SPARQL & Moyenne (2) \\
\hline
6 & Adaptabilité & Facilité d'adaptation au contexte burkinabè & Élevée (3) \\
\hline
7 & Traçabilité & Capacité à associer les connaissances aux sources & Élevée (3) \\
\hline
8 & Prise en compte des savoirs locaux & Intégration des noms vernaculaires, recettes multi-plantes & Élevée (3) \\
\hline
\end{tabular}
\end{table}

Nous avons attribué à chaque critère une pondération proportionnelle à son importance. Les critères fortement pondérés (3) étaient déterminants pour représenter le contexte burkinabè, tandis que les critères moyennement pondérés (2) autorisaient davantage de compromis. Nous avons calculé le score global en multipliant chaque note par la pondération correspondante, en additionnant les produits et en divisant cette somme par la somme des pondérations. Le résultat est rapporté sur 5.

\subsubsection{Notation des ontologies}

Nous avons noté chaque ontologie de 1 (très faible) à 5 (très élevé) pour chacun des critères. Le tableau \ref{tab:notation_ontologies} présente les résultats.

\begin{table}[H]
\centering
\caption{Notation des ontologies candidates}
\label{tab:notation_ontologies}
\begin{tabular}{|L{3.5cm}|C{0.8cm}|C{0.8cm}|C{0.8cm}|C{0.8cm}|C{0.8cm}|C{0.8cm}|C{0.8cm}|C{0.8cm}|C{1.5cm}|}
\hline
\textbf{Ontologie} & \textbf{C1} & \textbf{C2} & \textbf{C3} & \textbf{C4} & \textbf{C5} & \textbf{C6} & \textbf{C7} & \textbf{C8} & \textbf{Score} \\
\hline
WATRIMed & 4 & 4 & 2 & 3 & 5 & 4 & 3 & 5 & 3.7 \\
\hline
ontoMEDTRAD & 5 & 5 & 5 & 4 & 4 & 5 & 4 & 5 & 4.7 \\
\hline
ATHMed & 4 & 4 & 3 & 3 & 4 & 3 & 2 & 3 & 3.2 \\
\hline
Deep Ontologies & 3 & 4 & 2 & 2 & 3 & 3 & 2 & 4 & 2.9 \\
\hline
FLOPO / PO & 2 & 3 & 5 & 4 & 5 & 2 & 1 & 1 & 2.7 \\
\hline
\end{tabular}
\end{table}

Nous avons analysé les notes afin de comparer la couverture du domaine et la prise en compte des savoirs locaux par les ontologies africaines. Nous avons également examiné la disponibilité et l'interopérabilité des ontologies botaniques, ainsi que les différences entre WATRIMed et ontoMEDTRAD concernant la richesse du modèle et l'accessibilité de la documentation.

Nous avons accordé une attention particulière à la disponibilité (C3) et à la documentation (C4). Nous avons évalué les ressources effectivement accessibles. WATRIMed est décrite par deux articles de congrès, dont l'un compte deux pages, mais son fichier d'ontologie n'est pas téléchargeable. La thèse consacrée à ontoMEDTRAD fournit la liste des classes et des propriétés, les métriques ainsi que l'origine des concepts.

\subsubsection{Sélection de l'ontologie}

WATRIMed \cite{watrimed2019, watrimed_design_2019} pourrait se distinguer par la richesse de son modèle documenté, qui compte 472 concepts et 75 propriétés, ainsi que par son alignement sur l'ontologie fondationnelle BioTopLite2. Toutefois, les deux publications n'en décrivent que les classes supérieures et les relations principales, et le fichier de l'ontologie n'est pas disponible pour une réutilisation directe.

Nous avons retenu ontoMEDTRAD \cite{kouame_systeme_2018}, qui a obtenu le meilleur score. Cette ontologie couvre l'Afrique de l'Ouest, représente les recettes multi-plantes, leurs ingrédients, les parties utilisées et les modes de préparation, et documente l'origine de ses concepts à partir de quinze ressources terminologiques et ontologiques. Sa thèse est accessible. Nous avons aussi pris en compte ses limites : sa validation globale reste annoncée comme perspective et son jeu de données comprend 22 plantes et 31 remèdes antipaludéens \cite{kouame_visual_2020}.

L'ontologie retenue fournira le cadre conceptuel pour extraire et organiser les connaissances. Elle guidera la constitution du corpus et son adaptation au contexte burkinabè.

\textbf{Livrable} : Ontologie ontoMEDTRAD sélectionnée.

% ------------------------------------------------------------
% II.2.2. Phase 2 : Définition des critères de recherche et de sélection
% ------------------------------------------------------------

\subsection{Phase 2 : Définition des critères de recherche et de sélection des documents}
\label{subsec:phase2}

\textbf{Objectif} : Définir les critères qui guideront la recherche et la sélection des documents scientifiques.

Nous avons défini deux catégories de critères, l'une pour rechercher les informations et l'autre pour sélectionner les documents. Elles visaient à assurer la pertinence, la qualité et la cohérence des sources retenues.

\subsubsection{Critères de recherche}

Nous avons retenu les critères suivants pour guider la recherche :

\begin{itemize}
\item \textbf{Les plantes} : espèces végétales documentées au Burkina Faso, avec identification botanique (nom scientifique, famille, synonymes).
\item \textbf{Les parties utilisées} : feuilles, racines, écorces, fruits, graines, fleurs, tiges, etc.
\item \textbf{Les pathologies et usages médicinaux} : paludisme, infections respiratoires, digestives, diabète, hypertension, etc.
\item \textbf{Les modes de préparation} : décoction, macération, infusion, poudre, extrait alcoolique, etc.
\item \textbf{Les niveaux de preuve} : savoirs traditionnels, études in vitro, études in vivo, essais cliniques.
\item \textbf{Les propriétés pharmacologiques} : antimicrobiennes, antipaludiques, anti-inflammatoires, antioxydantes, toxicité, etc.
\item \textbf{La composition chimique} : alcaloïdes, flavonoïdes, tanins, saponines, huiles essentielles, etc.
\end{itemize}

Ils couvrent les principales dimensions des connaissances sur les plantes médicinales et ont permis de cibler les informations utiles à l'étude.

\subsubsection{Critères de sélection des documents}

Nous avons sélectionné les documents selon les critères suivants :

\begin{itemize}
\item \textbf{Pertinence} : le document traite des plantes médicinales du Burkina Faso.
\item \textbf{Identification de la plante} : le nom scientifique est précisé.
\item \textbf{Disponibilité d'une référence scientifique} : le document est publié dans une revue à comité de lecture, une thèse, un mémoire ou un rapport de recherche.
\item \textbf{Fiabilité} : les sources sont vérifiables et citées.
\item \textbf{Actualité} : les documents récents sont privilégiés, sans exclure les travaux fondateurs.
\end{itemize}

Ces critères nous ont permis d'écarter les documents non scientifiques ou non pertinents et de constituer un corpus cohérent pour l'étude.

\textbf{Livrable} : grille de critères de recherche et de sélection.

% ------------------------------------------------------------
% II.2.3. Phase 3 : Constitution du corpus scientifique
% ------------------------------------------------------------

\subsection{Phase 3 : Constitution du corpus scientifique}
\label{subsec:phase3}

\textbf{Objectif} : Rassembler un corpus documentaire fiable et pertinent sur les plantes médicinales du Burkina Faso.

Nous avons constitué le corpus en quatre étapes. Nous avons recherché les documents dans Google Scholar, PubMed et ResearchGate, puis consulté burkinaherbiermed.org et les données de l'INERA. Nous avons inclus les documents portant sur les plantes médicinales du Burkina Faso, précisant le nom scientifique des plantes et publiés ou validés par une institution scientifique avec des sources vérifiables. Nous avons écarté les documents non scientifiques, ceux sans identification botanique et les documents anciens sans valeur de référence. Nous avons classé les documents par type, plante, pathologie et région et les avons enregistrés dans un répertoire structuré. Le corpus comprend 35 documents scientifiques et les données de burkinaherbiermed.org et de l'INERA.

\textbf{Livrable} : un corpus de 35 documents scientifiques, complété par les données de burkinaherbiermed.org et de l'INERA.

% ------------------------------------------------------------
% II.2.4. Phase 4 : Adaptation de l'ontologie sélectionnée
% ------------------------------------------------------------

\subsection{Phase 4 : Adaptation de l'ontologie sélectionnée}
\label{subsec:phase4}

\textbf{Objectif} : Adapter le modèle conceptuel de l'ontologie sélectionnée au contexte burkinabè.

Nous avons examiné les concepts existants, ajouté des classes et des propriétés, modifié certaines relations et documenté les changements.

\subsubsection{Analyse des concepts existants}

Nous avons étudié la structure de l'ontologie ontoMEDTRAD telle qu'elle est décrite dans la thèse \cite{kouame_systeme_2018}. L'ontologie compte 144 concepts, 76 propriétés associatives, 16 propriétés de données, 789 individus et 2 734 axiomes. Elle est formalisée en logique de description ALCHOIQ, éditée sous Protégé~5 et vérifiée par les raisonneurs HermiT et Pellet.

Le fichier OWL de l'ontologie n'étant pas téléchargeable publiquement, nous avons reconstitué sa structure à partir de la thèse. Nous nous sommes appuyés sur les captures d'écran de Protégé (figures 5.15 à 5.20) et les tableaux récapitulatifs des concepts, propriétés et métriques (tableaux 5.5, 5.21 et 5.22). Cette reconstitution a servi de base à l'adaptation présentée au chapitre \ref{chap:conception_developpement}.

Nous avons reconstitué la structure conceptuelle de l'ontologie, ses classes, propriétés et relations, mais pas ses instances d'origine. Nous n'avons pas réutilisé les 789 individus d'ontoMEDTRAD, liés au contexte et aux sources des travaux d'origine. Cette phase a produit le modèle conceptuel adapté ; le peuplement par les individus extraits de la littérature scientifique est intervenu en phase 5.

Ses concepts principaux décrivent le parcours de soin, du praticien et du patient aux symptômes, aux maladies, aux ressources médicinales, à l'examen clinique, aux recettes, aux ingrédients, aux parties utilisées, aux préparations et aux remèdes. L'ontologie distingue cinq types de plantes et neuf parties utilisées. Elle comprend aussi une classe pour les références scientifiques associées aux ressources, aux recettes et aux remèdes. Les concepts secondaires couvrent les contre-indications, les propriétés thérapeutiques, la toxicité, l'étiologie et la classification des maladies selon la CIM-10. Les relations de base expriment la subsomption et la relation partie-tout. Soixante-seize propriétés associatives complètent le modèle, dont la hiérarchie \textit{aPourForme}, utilisée pour l'iconisation des plantes.

Nous avons ensuite évalué ce modèle au regard du contexte burkinabè. Il couvre les concepts de base et les raisonneurs testés n'y ont signalé aucune incohérence, mais il nécessitait des adaptations pour représenter les noms vernaculaires en langues locales, les maladies spécifiques, les niveaux de preuve, les dosages et la traçabilité détaillée des sources. Sa validation globale reste une perspective, et son jeu de données se limite à 22 plantes et 31 remèdes antipaludéens.

\subsubsection{Ajout de classes}

La comparaison des fichiers du modèle conceptuel adapté et du schéma conservé dénombre 40 classes ajoutées. Le modèle adapté déclare également 25 propriétés objet au total, dont 20 absentes du schéma conservé, et 20 propriétés de données, dont 15 ajoutées. Les classes ajoutées couvrent notamment les noms vernaculaires, les usages thérapeutiques, les niveaux de preuve, les sources, les préparations, des parties végétales complémentaires, les composés chimiques, les activités pharmacologiques et les effets toxicologiques. Ces ajouts définissent le modèle conceptuel ; les individus et données extraits sont intégrés lors de la phase suivante.

Ces ajouts structurent des connaissances locales et scientifiques à différents niveaux. Ils ne préchargent pas les valeurs, comme les noms, pathologies, régions ou composés ; celles-ci proviennent des sources du corpus et sont enregistrées dans la base avec leur provenance.

\subsubsection{Ajout de propriétés}

Nous avons ajouté des propriétés de données pour les noms locaux, le niveau de preuve, le type d'étude, les références, le dosage et la toxicité, en complément de propriétés textuelles pour les noms vernaculaires et les descriptions d'activités ou d'effets. Des propriétés objet relient également l'usage à sa maladie cible, à la partie végétale utilisée, à la préparation et à la source ; d'autres décrivent les associations avec des composés, activités pharmacologiques, risques et localités. Les noms techniques exacts sont ceux déclarés dans la version courante du fichier Turtle.

L'enrichissement ainsi obtenu saisit des informations fines (noms vernaculaires, coordonnées, niveaux de preuve, références précises) que l'ontologie d'origine ne documentait pas pour le Burkina Faso.

\subsubsection{Modification de relations}

Nous avons structuré les liens entre plante, usage, maladie, partie utilisée et préparation, puis généralisé la traçabilité documentaire. La propriété objet \textit{estDocumenteePar} relie les ressources du graphe à une source scientifique; des sous-propriétés spécialisées permettent de préciser si la source documente une plante, un usage, un nom vernaculaire, une activité, un composé ou un effet toxicologique. \textit{aPourNiveauDePreuve} rattache un usage au niveau déclaré par la source. Ces liens ne signifient pas que les usages ou activités constituent des recommandations thérapeutiques ni une validation clinique.

Ces modifications assurent la traçabilité des connaissances et permettent d'exprimer des relations que le modèle d'origine, fondé sur la relation « est un » et sur la relation partie-tout, ne couvrait pas.

\subsubsection{Documentation des adaptations}

Nous avons consigné toutes les modifications dans un journal : chaque ajout est justifié, chaque choix est explicité, et les classes, propriétés et relations ajoutées, modifiées ou supprimées y sont décrites une à une.

Cette documentation n'est pas une formalité. Elle rend la démarche reproductible par un tiers et facilite les évolutions ultérieures de l'ontologie.

\textbf{Livrable} : Le modèle conceptuel adapté de l'ontologie, sans les données instanciées.

% ------------------------------------------------------------
% II.2.5. Phase 5 : Enrichissement de l'ontologie
% ------------------------------------------------------------

\subsection{Phase 5 : Enrichissement des connaissances à partir de la littérature scientifique}
\label{subsec:phase5}

\textbf{Objectif} : Extraire les connaissances du corpus pour alimenter la base organisée selon l'ontologie.

À partir du modèle conceptuel adapté en phase 4, nous avons enrichi l'ontologie avec des individus et des données extraits des sources scientifiques. Nous avons préparé ces données structurées et sourcées pour leur exploitation dans la base orientée graphe Neo4j.

Nous avons conçu les consignes d'extraction, procédé à l'extraction assistée par modèle de langue, contrôlé les données obtenues, puis structuré et préparé ces données pour leur intégration dans la base.

\subsubsection{Conception des prompts d'extraction}

Nos consignes d'extraction ciblaient les noms scientifiques et vernaculaires, la famille botanique, les parties utilisées, les usages, les préparations, les activités pharmacologiques, le niveau de preuve et la source. Le chapitre \ref{chap:conception_developpement} présente un exemple de consigne, avec le format JSON et les champs demandés.

Nous avons employé une même consigne pour les documents traités. Les champs absents reçoivent la valeur \texttt{null} ou une liste vide, et aucune valeur n'est complétée par déduction.

\subsubsection{Exécution de l'extraction}

Nous avons soumis les documents à gpt-6-luna dans l'environnement OpenWork, un par un, et enregistré les résultats en JSON. L'extraction a porté sur dix des 35 documents du corpus, principalement des enquêtes ethnobotaniques ; les 25 autres n'ont pas été extraits.

\subsubsection{Contrôle de la qualité des extractions}

Nous avons repris le niveau de preuve indiqué par la source lorsqu'il était explicite, qu'il s'agisse d'un savoir traditionnel, d'une étude \textit{in vitro}, d'une étude \textit{in vivo} ou d'un essai clinique. Nous n'avons pas évalué indépendamment la solidité des faits et avons laissé le champ vide lorsque la source ne précisait pas ce niveau. Un expert externe n'a pas relu les extractions. Nous avons contrôlé leur structure, puis comparé une partie des faits aux documents d'origine à la phase 8.

\subsubsection{Structuration des données}

Nous avons organisé les données selon le modèle de l'ontologie dans un classeur d'individus, par type de connaissance (sources, plantes, noms vernaculaires, usages, préparations, maladies, activités, composés, toxicologie et localités), puis les avons transformées en individus au format Turtle.

\subsubsection{Préparation de l'intégration dans la base de connaissances}

Le fichier OntoPlantMedBF-peuplee.ttl contient des individus extraits de dix documents scientifiques, en plus du schéma adapté. Ce fichier constitue l'ontologie enrichie importée dans Neo4j à la phase 6. Les données de burkinaherbiermed.org et de l'INERA, consultées lors de la constitution du corpus, n'y ont pas été intégrées.

\textbf{Livrable} : OntoPlantMedBF, ontologie enrichie avec les données extraites et leurs sources.

% ------------------------------------------------------------
% II.2.6. Phase 6 : Construction et interrogation de la base
% ------------------------------------------------------------

\subsection{Phase 6 : Construction et interrogation de la base de connaissances}
\label{subsec:phase6}

\textbf{Objectif} : Importer et exploiter dans Neo4j l'ontologie réajustée \textit{OntoPlantMedBF}.

Nous avons importé dans Neo4j l'ontologie réajustée \textit{OntoPlantMedBF} et ses données, en les représentant sous forme de nœuds, propriétés et relations du graphe. Nous avons vérifié la correspondance avec le modèle ontologique et testé les requêtes de la base.

\subsubsection{Import et ajustement du modèle dans le graphe}

Nous avons importé les classes et propriétés d'OntoPlantMedBF dans Neo4j sous forme de labels, propriétés et relations. Les principaux nœuds représentent les plantes, maladies, usages, parties, préparations, sources, composés chimiques et régions. Les relations du graphe reprennent les liens du modèle, avec des noms adaptés aux conventions Neo4j, et les propriétés conservent notamment les noms, niveaux de preuve, DOI et coordonnées géographiques.

Nous avons ainsi importé dans le graphe la structure réajustée d'OntoPlantMedBF et ses données. Les entités sont représentées par des nœuds, leurs liens par des relations et leurs caractéristiques par des propriétés, conformément au modèle ontologique.

\subsubsection{Import des données}

Nous avons chargé les données préparées, créé ou rapproché les nœuds et relations au moyen de requêtes Cypher, puis vérifié l'intégrité et la cohérence de l'import avec le modèle ontologique.

Après l'import, nous avons contrôlé que les données, les relations et les propriétés avaient été correctement reprises.

\subsubsection{Requêtes de test}

Nous avons testé quatre familles de requêtes : les requêtes simples sur les plantes associées à une maladie, les requêtes complexes sur les recettes multi-plantes, les requêtes de traçabilité et les requêtes spatiales par région.

Ces vérifications nous ont permis de contrôler le fonctionnement de la base et les résultats renvoyés par les requêtes.

\textbf{Livrable} : Une base Neo4j peuplée et interrogeable, complémentaire d'OntoPlantMedBF.

% ------------------------------------------------------------
% II.2.7. Phase 7 : Exploitation par l'IA et développement web
% ------------------------------------------------------------

\subsection{Phase 7 : Exploitation par IA et développement de l'application web}
\label{subsec:phase7}

\textbf{Objectif} : Développer le système d'accès intelligent et l'interface utilisateur.

Nous avons mis en place le système RAG en définissant son architecture, en construisant la chaîne de traitement et en intégrant ses composants.

\subsubsection{Architecture du système}

L'architecture relie l'utilisateur à Streamlit, LangChain, Neo4j et Llama 3 jusqu'à la réponse. Elle associe la recherche dans la base à la génération par le modèle de langue et conserve le lien entre les informations et leurs sources. Le chapitre \ref{chap:conception_developpement} décrit la chaîne de traitement, de l'interprétation de la question à l'affichage des références.

\subsubsection{Intégration des composants}

Nous avons connecté l'application à Neo4j et au modèle de langue, puis configuré les consignes de génération des réponses. Cette intégration a réuni les composants du système RAG.

\subsubsection{Développement de l'application web}

Nous avons conçu l'application autour d'une page d'accueil avec chatbot, d'une recherche par plante, d'une recherche par pathologie et de fiches descriptives. Nous y avons intégré le système RAG et la gestion de l'historique. L'interface permet d'interroger la base en langage naturel et de recevoir des réponses reliées à leurs sources. Le chapitre \ref{chap:conception_developpement} en décrit la conception et les composants.

\subsubsection{Tests et validation}

Nous avons mené des tests fonctionnels pour vérifier le fonctionnement de l'application et des tests utilisateurs pour examiner sa facilité d'utilisation et la pertinence des réponses.

\textbf{Livrable} : Un système RAG fonctionnel et une application web fonctionnelle.

% ------------------------------------------------------------
% II.2.8. Phase 8 : Évaluation et validation
% ------------------------------------------------------------

\subsection{Phase 8 : Évaluation et validation}
\label{subsec:phase8}

\textbf{Objectif} : Évaluer la qualité de l'ontologie, du système RAG et de l'application web.

Nous avons évalué trois composantes : l'ontologie, le système RAG et l'application web.

\subsubsection{Évaluation de l'ontologie}

Nous avons évalué l'ontologie selon les quatre critères présentés dans le tableau \ref{tab:eval_ontologie}.

\begin{table}[H]
\centering
\caption{Critères d'évaluation de l'ontologie}
\label{tab:eval_ontologie}
\begin{tabular}{|L{3cm}|L{6cm}|L{4cm}|}
\hline
\textbf{Critère} & \textbf{Description} & \textbf{Méthode} \\
\hline
Couverture & Capacité à représenter les connaissances du corpus & Vérification manuelle \\
\hline
Cohérence & Absence de contradictions & Requêtes de test \\
\hline
Traçabilité & Association des connaissances aux sources & Vérification des relations \\
\hline
Extensibilité & Facilité d'ajout de nouvelles connaissances & Test d'ajout \\
\hline
\end{tabular}
\end{table}

\subsubsection{Évaluation du système RAG}

Nous avons évalué le système RAG selon les quatre critères présentés dans le tableau \ref{tab:eval_rag}.

\begin{table}[H]
\centering
\caption{Critères d'évaluation du système RAG}
\label{tab:eval_rag}
\begin{tabular}{|L{3cm}|L{6cm}|L{4cm}|}
\hline
\textbf{Critère} & \textbf{Description} & \textbf{Métrique} \\
\hline
Pertinence & Les réponses répondent à la question & Précision \\
\hline
Fidélité & Les réponses sont fidèles aux sources & Rappel \\
\hline
Exactitude & Les réponses sont correctes & F1-score \\
\hline
Traçabilité & Les sources sont affichées & Taux de couverture \\
\hline
\end{tabular}
\end{table}

Nous avons constitué un jeu de 50 questions de référence portant sur les plantes, les usages et leurs sources, puis comparé les réponses du système aux réponses attendues.

\subsubsection{Évaluation de l'application web}

Nous avons évalué l'application web selon les quatre critères présentés dans le tableau \ref{tab:eval_web}.

\begin{table}[H]
\centering
\caption{Critères d'évaluation de l'application web}
\label{tab:eval_web}
\begin{tabular}{|L{3cm}|L{6cm}|L{4cm}|}
\hline
\textbf{Critère} & \textbf{Description} & \textbf{Méthode} \\
\hline
Facilité d'utilisation & L'interface est intuitive & Questionnaire \\
\hline
Compréhension & Les réponses sont claires & Questionnaire \\
\hline
Utilité & L'application répond aux besoins & Questionnaire \\
\hline
Performance & Temps de réponse & Mesure \\
\hline
\end{tabular}
\end{table}

Nous avons sollicité des étudiants, des chercheurs et des tradipraticiens pour cette évaluation.

\textbf{Livrable} : Rapport d'évaluation et de validation.

% ============================================================
% II.3. SYNTHÈSE ET TRANSITION
% ============================================================

\section{Synthèse et transition}
\label{sec:synthese}

\subsection{Tableau de correspondance Phases - Livrables}

Le tableau \ref{tab:correspondance_phases} récapitule les objectifs et les livrables des huit phases.

\begin{table}[H]
\centering
\caption{Correspondance entre les phases, les objectifs et les livrables}
\label{tab:correspondance_phases}
\begin{tabular}{|L{2.2cm}|L{4.5cm}|L{6.8cm}|}
\hline
\textbf{Phase} & \textbf{Objectif} & \textbf{Livrable} \\
\hline
Phase 1 & Sélectionner l'ontologie la plus pertinente & Ontologie ontoMEDTRAD sélectionnée \\
\hline
Phase 2 & Définir les critères de recherche et de sélection & Grille de critères \\
\hline
Phase 3 & Rassembler un corpus fiable et pertinent & Corpus de 35 documents \\
\hline
Phase 4 & Adapter le modèle conceptuel de l'ontologie & Modèle conceptuel adapté \\
\hline
Phase 5 & Extraire les connaissances et enrichir l'ontologie & Ontologie enrichie en données \\
\hline
Phase 6 & Importer OntoPlantMedBF dans la base graphe & OntoPlantMedBF intégrée à Neo4j, base interrogeable \\
\hline
Phase 7 & Développer le système RAG et l'application web & Système RAG et application web \\
\hline
Phase 8 & Évaluer la qualité du système & Rapport d'évaluation \\
\hline
\end{tabular}
\end{table}

% ============================================================
% CONCLUSION DU CHAPITRE
% ============================================================

\section*{Conclusion du chapitre}
\addcontentsline{toc}{section}{Conclusion du chapitre}

Ce chapitre a exposé le matériel et la méthodologie de notre recherche. Le matériel réunit un corpus de 35 documents scientifiques, les données de burkinaherbiermed.org et de l'INERA, cinq ontologies candidates et un environnement technique bâti sur gpt-6-luna pour l'extraction, Llama 3 8B sur Google Colab pour le RAG, Neo4j Desktop, LangChain et Streamlit.

Nous avons suivi huit phases : comparaison des ontologies, définition des critères de recherche et de sélection, constitution du corpus, adaptation de l'ontologie, extraction des connaissances, construction et interrogation de la base, développement du système d'accès et de l'application, puis évaluation. Nous avons décrit les actions menées et les livrables associés.

Nous avons choisi l'ontologie avant de constituer le corpus afin d'assurer la cohérence entre le modèle conceptuel et les données collectées, et de cibler les informations pertinentes dès la collecte.

Nous signalons deux limites. Les niveaux de preuve associés aux connaissances proviennent des sources et n'ont pas fait l'objet d'une évaluation indépendante. De plus, un expert externe n'a pas validé les extractions. Ces limites sont à prendre en compte dans les travaux ultérieurs. Nous avons néanmoins conservé le lien entre chaque information et sa source.

Le chapitre \ref{chap:resultats_discussion} présente les résultats de l'évaluation menée en phase 8.

Le chapitre suivant présente la conception et le développement du système, notamment l'ontologie adaptée, la base de connaissances, le système RAG et l'application web.
